% !TEX program = xelatex
\documentclass[11pt,openany,fontset=fandol]{ctexbook}
\usepackage{geometry,amsmath,amssymb,amsthm,bm,mathtools,booktabs,array,enumitem,float}
\usepackage{tikz,fancyhdr,hyperref}
\usetikzlibrary{arrows.meta,calc,patterns,angles,quotes,decorations.markings,decorations.pathmorphing}
\theoremstyle{definition}\newtheorem{example}{例}[chapter]\newtheorem{drill}{练习}[chapter]
\newenvironment{solution}{\par\smallskip\noindent\textbf{解}\quad\ignorespaces}{\par\medskip}
\newenvironment{drillsolution}[1]{\subsubsection*{练习\ref{#1}详解}\phantomsection\label{sol:#1}\ignorespaces}{\par\medskip}
\newenvironment{question}{\par\medskip\noindent\begin{minipage}{\linewidth}}{\end{minipage}\par\medskip}
\newcommand{\dd}{\mathop{}\!\mathrm d}
\newcommand{\e}{\mathrm e}
\newcommand{\fig}[3]{\begin{figure}[H]\centering #3\caption{#2}\label{#1}\end{figure}}
\newcommand{\answerlink}[1]{\par\smallskip\noindent\hyperref[sol:#1]{详解}}
\tikzset{axis/.style={->,>=Stealth,line width=.45pt},world/.style={line width=.8pt},flow/.style={->,>=Stealth,line width=.7pt},guide/.style={dashed,line width=.35pt},dim/.style={<->,>=Stealth,line width=.35pt},every node/.style={font=\small},force/.style={->,>=Stealth,very thick},spring/.style={decorate,decoration={zigzag,segment length=3pt,amplitude=2pt}}}
% Common visual style for all first-round physics lecture projects.
% Each project includes an identical local copy of this file.
\geometry{paperwidth=176mm,paperheight=250mm,inner=21mm,outer=18mm,top=20mm,bottom=20mm,headheight=14pt,headsep=6mm,footskip=10mm}
\hypersetup{hidelinks}
\setlength{\parindent}{2em}\setlength{\parskip}{3pt}\setlength{\emergencystretch}{2em}
\linespread{1.18}\raggedbottom\setcounter{tocdepth}{1}
\widowpenalty=10000\clubpenalty=10000\displaywidowpenalty=10000
\setlength{\intextsep}{8pt plus1pt minus1pt}
\makeatletter
\g@addto@macro\normalsize{\setlength{\abovedisplayskip}{8pt plus2pt minus2pt}\setlength{\belowdisplayskip}{8pt plus2pt minus2pt}\setlength{\abovedisplayshortskip}{0pt plus2pt}\setlength{\belowdisplayshortskip}{6pt plus2pt minus2pt}}
\makeatother
\providecommand{\Needspace}[1]{}
\newdimen\LectureAvailableSpace
\renewcommand{\Needspace}[1]{\par\penalty-100\begingroup\LectureAvailableSpace=\pagegoal\advance\LectureAvailableSpace by-\pagetotal\ifdim\LectureAvailableSpace<#1\relax\ifdim\pagetotal>0pt\newpage\fi\fi\endgroup}
\renewenvironment{solution}{\par\smallskip\Needspace{3\baselineskip}\noindent\textbf{解}\quad\ignorespaces}{\par\medskip}
\ctexset{chapter={name={第,讲},number=\chinese{chapter},format=\LARGE\bfseries\centering,beforeskip=16pt,afterskip=22pt},section={format=\Large\bfseries},subsection={format=\large\bfseries},subsubsection={format=\normalsize\bfseries,beforeskip=10pt,afterskip=4pt}}
\setlist[enumerate]{leftmargin=2em,itemsep=2pt,topsep=3pt}
\pagestyle{fancy}\fancyhf{}\fancyhead[LE,RO]{\small\thepage}\fancyhead[LO]{\small 物理竞赛一轮讲义}
\renewcommand{\headrulewidth}{0pt}
\fancypagestyle{plain}{\fancyhf{}\fancyfoot[C]{\small\thepage}\renewcommand{\headrulewidth}{0pt}}
\renewenvironment{drillsolution}[1]{\par\Needspace{4\baselineskip}\subsubsection*{练习\ref{#1}详解}\phantomsection\label{sol:#1}\nobreak\noindent\textbf{解}\quad\ignorespaces}{\par\medskip}
\ifdefined\exercisesolution
\renewenvironment{exercisesolution}[1]{\par\Needspace{4\baselineskip}\subsubsection*{习题\ref{#1}详解}\phantomsection\label{sol:#1}\nobreak\noindent\textbf{解}\quad\ignorespaces}{\par\medskip}
\fi
\newcommand{\LectureAnswerSection}{\par\section*{练习解答}\addcontentsline{toc}{section}{练习解答}}
\newcommand{\LectureHeader}[2]{\fancyhead[RE]{\small #1\quad #2}}
\newcommand{\LectureCover}[3]{\begin{titlepage}\centering\vspace*{32mm}{\Huge\bfseries #1\par}\vspace{12mm}{\Large 物理竞赛一轮讲义\par}\vspace{18mm}{\LARGE #2\\[5mm]#3\par}\vfill{\large 2026\par}\end{titlepage}}

\LectureHeader{第四讲}{动力学过程与综合}
\begin{document}\frontmatter\LectureCover{牛顿运动定律}{第四讲}{动力学过程与综合}\tableofcontents\mainmatter
\setcounter{chapter}{3}\chapter{动力学过程与综合}
\section{以事件连接动力学过程}
\subsection{先确定阶段边界，再使用运动学}
动力学综合题常由多个简单过程组成，困难不在某条单独公式，而在阶段边界的识别。撤力、换接触面、速度变零、两物体速度相同、碰撞以及支持力降为零，都会改变受力或约束。先写每段实际受力，再判断其加速度是否恒定，最后以边界事件连接相邻两段，才能避免将全程误当作一段匀加速运动。

若第 $j$ 段从时刻 $t_j$ 开始，恒加速度为 $a_j$，应写 $v=v_j+a_j(t-t_j)$、$x=x_j+v_j(t-t_j)+a_j(t-t_j)^2/2$。式中的局部时间与绝对时刻不同。非冲击切换通常保持位置、速度连续；加速度可以跳变。碰撞则可能在很短时间内使速度改变，但碰前、碰后的位置应相容，不能让两物体穿越。

\fig{n4t-events}{连续速度与分段加速度分别承担不同的衔接条件}{\begin{tikzpicture}
\draw[axis](0,0)--(5.4,0)node[right]{$t$};\draw[axis](0,0)--(0,2.4)node[above]{$v$};\draw[world](0,0)--(1.8,1.8)--(3,1.2)--(5,1.2);\draw[guide](1.8,0)--(1.8,2.2) (3,0)--(3,2.2);\node[below]at(1.8,0){撤力};\node[below]at(3,0){共同运动};\node at(4,1.7){$a=0$};
\end{tikzpicture}}
图\ref{n4t-events}在折点处速度连续，斜率改变。假如物块与长板速度相同，不能立刻预设以后共同运动；应先算共同加速度所需静摩擦，若未超上限才能转入共同运动，否则相对滑动会换向或继续。

\fig{n4deep:match}{分段计算在同一边界继承位置与速度，斜率可以改变}{\begin{tikzpicture}
\begin{scope}\draw[axis](0,0)--(4.4,0)node[right]{$t$};\draw[axis](0,0)--(0,3.4)node[above]{$x$};\draw[world,domain=0:2,samples=30]plot(\x,{.4*(\x)^2});\draw[world,blue,domain=2:3.6,samples=30]plot(\x,{1.6+1.6*(\x-2)-.4*(\x-2)^2});\draw[guide](2,0)--(2,1.6);\fill(2,1.6)circle(2pt)node[above left]{$x_j$};\draw[red,dashed](1.3,.48)--(2.7,2.72);\node[below]at(2,0){$t_j$};\end{scope}
\begin{scope}[xshift=5.5cm]\draw[axis](0,0)--(4.1,0)node[right]{$t$};\draw[axis](0,0)--(0,2.4)node[above]{$v$};\draw[world](0,0)--(2,1.6);\draw[world,blue](2,1.6)--(3.6,.32);\fill(2,1.6)circle(2pt)node[above]{$v_j$};\draw[guide](2,0)--(2,1.6);\node[below]at(2,0){$t_j$};\node at(2,-.7){新一段用 $\tau=t-t_j$};\end{scope}
\end{tikzpicture}}
图\ref{n4deep:match}左图在边界处共享位置和切线，右图在同一时刻共享速度值，但两边斜率不同。后一段的初值应当是 $(x_j,v_j)$，即 $x=x_j+v_j\tau+\tfrac12a_j\tau^2$。若误把绝对时刻 $t$ 直接当作新一段历时，边界处便通常出现位置或速度的虚假跳跃。图中的连续性以切换不产生冲量为前提，碰撞须另用碰前、碰后速度连接。

\subsection{阻力由速度方向决定}
大小恒定的滑动阻力或空气阻力，其方向通常与当时速度相反。飞行器先上升再下落时，同一个大小 $f$ 在竖直方程中的符号需要改变；恢复升力但还在向下运动的阶段，升力与阻力可能共同向上制动。不能因为“恢复了动力”就直接套用先前上升阶段的加速度。

速度恰为零的时刻是阶段边界，恒大小阻力模型在这一点本身可能不光滑。判断接下来向哪边运动，要用外力和允许的静态条件；在采用纯运动阻力近似的飞行器模型中，可分别求上升终点前、下降起点后两个极限。制动距离还须计入恢复动力前已经损失的高度，而不是只看恢复时是否仍在空中。

\fig{n4deep:forceswitch}{阻力随速度改变方向，动力恢复时还要保留当时速度方向}{\begin{tikzpicture}
\foreach \x/\title in {0/无动力上升,3.6/无动力下降,7.2/恢复动力制动}{\fill(\x,0)circle(2pt);\draw[force,red](\x,0)--(\x,-1.2)node[below]{$mg$};\node at(\x,-2){\title};}
\draw[force,blue](.4,0)--(.4,-.7)node[right]{$f$};\draw[flow](-.6,0)--(-.6,1)node[above]{$v>0$};
\draw[force,blue](3.6,0)--(3.6,.7)node[above]{$f$};\draw[flow](3,0)--(3,-1)node[left]{$v<0$};
\draw[force,blue](7.2,0)--(7.2,.7)node[above]{$f$};\draw[force](7.7,0)--(7.7,1.5)node[above]{$F$};\draw[flow](6.6,0)--(6.6,-1)node[left]{$v<0$};
\end{tikzpicture}}
图\ref{n4deep:forceswitch}三图统一向上取正，分别对应 $ma=-mg-f$、$ma=-mg+f$ 和 $ma=F-mg+f$。第三图速度向下而加速度可以向上，表示正在制动；必须把制动距离计入总下降高度。待速度真正降到零并重新向上后，阻力才再次向下。由此可以逐个事件决定方程，而不必记忆彼此容易混淆的加速度。

地面摩擦中的停止事件还可能引入一段停留。例如物体原来向右滑动，恒定右向驱动力为 $P<\mu_kmg$，滑动阶段向左减速。到零速度后，若 $P\le f_{s,\max}$，地面静摩擦即可平衡驱动力，物体保持静止，不能继续按 $P-\mu_kmg<0$ 向左加速。只有另一个作用使它真正向左滑动时，动摩擦才改为向右。重复碰撞题中某个滑块先停下后等待追撞，使用的就是这种静止分支。

\fig{n4deep:stop}{停车后静摩擦维持静止，原减速直线不再适用}{\begin{tikzpicture}
\begin{scope}\node[draw]at(0,0){$m$};\draw[force](.4,0)--(1.3,0)node[above]{$P$};\draw[force,blue](-.4,0)--(-1.7,0)node[above]{$\mu_kmg$};\draw[flow](-.4,.8)--(.8,.8)node[right]{$v>0$};\node at(0,-.8){滑动阶段};\end{scope}
\begin{scope}[xshift=4cm]\draw[axis](0,0)--(4,0)node[right]{$t$};\draw[axis](0,-1)--(0,2.3)node[above]{$v$};\draw[world](0,1.8)--(2.4,0)--(3.6,0);\draw[red,dashed](2.4,0)--(3.5,-.825);\node[below left]at(2.4,0){$t_s$};\node at(2,2){停止后 $f_s=P$};\node[red]at(2.8,-1.3){虚线不满足摩擦方向};\end{scope}
\end{tikzpicture}}
图\ref{n4deep:stop}中 $0<P<\mu_kmg$，停车前 $a=(P-\mu_kmg)/m<0$；若 $P\le f_{s,\max}$，停车后取静摩擦 $f_s=P$、$a=0$。把红虚线继续到 $v<0$ 却仍保留向左动摩擦，会使摩擦沿运动方向做正功，直接暴露了阶段延续的错误。

\section{运动图像与受力的相互转换}
\subsection{斜率、面积与量纲}
在 $v$--$t$ 图中，斜率为加速度，乘质量得到合力；曲线下的有符号面积为位移。$a$--$t$ 图的面积则是速度变化，不是位移。$F$--$t$ 图的面积给冲量，$F$--$x$ 图的面积给功。它们纵横坐标不同，不能只凭“面积”二字互相代用。

如果给出稀疏的速度测量，先用确实属于同一阶段的相邻数据求斜率。斜面到粗糙水平面的交界时刻不一定正好在采样点，应由两段速度函数的连续性交点求出。再判断所问时刻落在哪一段，才能代入正确的速度函数。反求的摩擦因数还应使用该段实际支持力，而非无条件采用 $mg$。

有限采样只给区间平均加速度，不能单凭两个点就证明整个区间加速度恒定。讲义中的分段直线模型还使用了“光滑固定斜面”“恒定动摩擦”“平滑转接”等受力假设。实测数据若偏离直线，应先判断测量误差、交界区间和阻力模型，而不是选择部分点强行代入匀变速公式。根据图像得到的参数，应回代其余采样点检查是否同时满足。

\fig{n4t-graphs}{速度图的斜率与加速度图的高度逐段一致}{\begin{tikzpicture}
\begin{scope}\draw[axis](0,0)--(3.2,0)node[right]{$t$};\draw[axis](0,0)--(0,2.5)node[above]{$v$};\fill[blue!8](0,0)--(1,2)--(2.5,.5)--(2.5,0)--cycle;\draw[world](0,0)--(1,2)--(2.5,.5);\draw[guide](1,0)--(1,2);\node[below]at(1,0){$1$};\node[left]at(0,2){$2$};\node at(1.4,-.7){面积为 $\Delta x$};\end{scope}
\begin{scope}[xshift=5cm]\draw[axis](0,0)--(3.2,0)node[right]{$t$};\draw[axis](0,-1.3)--(0,2.5)node[above]{$a$};\fill[blue!8](0,0)rectangle(1,2);\fill[red!8](1,0)rectangle(2.5,-1);\draw[world](0,2)--(1,2) (1,-1)--(2.5,-1);\draw[guide](1,-1)--(1,2);\node[left]at(0,2){$2$};\node[left]at(0,-1){$-1$};\node[below right]at(1,0){$1$};\node at(1.4,-1.8){有符号面积为 $\Delta v$};\end{scope}
\end{tikzpicture}}
图\ref{n4t-graphs}采用一致单位：前一段速度从 $0$ 增到 $2$，历时 $1$，所以右图高度为 $2$；后一段速度在 $1.5$ 的时间内下降 $1.5$，右图高度为 $-1$。右图两块有符号面积之和为 $2-1.5=0.5$，正好等于左图的末速度减初速度。负面积使速度减小，而不是使速度立即变负。求两物体相对位移时，应积分速度差；求相对速度时，应积分加速度差。在图像中涂出两条曲线之间的区域，还须说明区域属于哪一时间区间，以及被积差是否发生变号。

\subsection{由速度反求的是合力}
测得一维有符号速度 $v(t)$，在惯性系内定质量模型给 $F_{\rm net}(t)=m\dot v(t)$。同一个合力可以由多种外力组合产生，因此仅靠速度函数不能唯一确定某一根绳张力、某一个驱动力或摩擦系数，除非其他力已知。若只给速率而运动方向不断变化，还不能只求速率导数就得到合力矢量，需要另外给出轨迹曲率与方向。

当不同物体的外力分别随时间变化，接触阶段应先取整体求共同加速度，再隔离求内部压力。压力到零时通常是分离候选点；分离后各自加速度按各自外力决定，之后两者间的速度差从零开始累积，而不是在分离瞬间就有有限的相对速度。

\fig{n4t-speed}{相同速率图不唯一决定合力，轨迹方向也提供信息}{\begin{tikzpicture}
\draw[world](-3,0)--(-.7,0);\draw[flow,blue](-2.6,0)--(-1.3,0)node[above]{$v$};\node at(-1.9,-.8){直线：$F=0$};\draw[world](1.7,0)circle(.9);\fill(2.6,0)circle(2pt);\draw[flow,blue](2.6,0)--(2.6,1.2)node[above]{$v$};\draw[force,red](2.6,0)--(1.7,0)node[midway,below]{$F$};\node at(1.7,-1.4){圆周：$F=mv^2/R$};
\end{tikzpicture}}
图\ref{n4t-speed}两物体速率都为常量，但左图速度方向不变，合力为零；右图切向速度与径向合力垂直，合力大小为 $mv^2/R$。只给速率图便无法区分这两种受力，所以逆向求力前必须确认所测量的是固定轴的速度分量，还是没有方向信息的速率。

\section{选择合适的微分变量}
\subsection{位置相关力与时间相关力}
若合力写成 $F(t)$，从初态积分有
$m\int_{v_0}^{v}\dd w=\int_{t_0}^{t}F(t')\dd t'$，
即 $m(v-v_0)=\int_{t_0}^{t}F(t')\dd t'$。
若力写成 $F(x)$，则按链式法则
$\dd v/\dd t=(\dd v/\dd x)(\dd x/\dd t)=v\dd v/\dd x$，
因而
$m\int_{v_0}^{v}w\dd w=\int_{x_0}^{x}F_{\rm net}(\xi)\dd\xi$。
左边用原函数 $mw^2/2$ 代入上下限，得到
\begin{equation}\label{n4eq-work}
\frac m2(v^2-v_0^2)=\int_{x_0}^xF_{\rm net}(\xi)\dd\xi.
\end{equation}
式中积分沿实际单调运动段进行。若先向右再向左，位移的符号应随路径改变，不能把路程一律当成正的 $\dd x$。在分段恒力近似下，式\eqref{n4eq-work}也可以从各段的 $v_{j+1}^2-v_j^2=2F_j\Delta x_j/m$ 相加直接得到。

平均力的概念因此必须附带权重。等时间段作用力的算术平均与冲量相联系；等位移段作用力的算术平均与功相联系。若作用力随位置线性变化，空间平均值等于两端平均；随时间线性变化时，这一平均并不自动等于空间平均。判断某个“平均值”能否代入前，先看要求累积的是速度还是速度平方。

\fig{n4deep:weights}{横轴决定累积量，同样的两级力不一定具有相同平均值}{\begin{tikzpicture}
\begin{scope}\draw[axis](0,0)--(3.6,0)node[right]{$t$};\draw[axis](0,0)--(0,2.6)node[above]{$F$};\fill[blue!8](0,0)rectangle(1,1);\fill[red!8](1,0)rectangle(3,2);\draw[world](0,1)--(1,1) (1,2)--(3,2);\node[left]at(0,1){$F_1$};\node[left]at(0,2){$F_2$};\node at(.5,.4){$\Delta t_1$};\node at(2,.4){$\Delta t_2$};\node at(1.5,-.7){总面积 $J$};\end{scope}
\begin{scope}[xshift=5cm]\draw[axis](0,0)--(3.6,0)node[right]{$x$};\draw[axis](0,0)--(0,2.6)node[above]{$F$};\fill[blue!8](0,0)rectangle(1.5,1);\fill[red!8](1.5,0)rectangle(3,2);\draw[world](0,1)--(1.5,1) (1.5,2)--(3,2);\node[left]at(0,1){$F_1$};\node[left]at(0,2){$F_2$};\node at(.75,.4){$\Delta x$};\node at(2.25,.4){$\Delta x$};\node at(1.5,-.7){总面积 $W$};\end{scope}
\end{tikzpicture}}
图\ref{n4deep:weights}对比两种给定方式。左图冲量为 $J=F_1\Delta t_1+F_2\Delta t_2$，时间平均力为 $J/(\Delta t_1+\Delta t_2)$；右图功为 $W=(F_1+F_2)\Delta x$，空间平均力才是 $(F_1+F_2)/2$。这两图不假定描述同一运动；它们强调，只有权重相等时才可直接取算术平均。

\subsection{线性阻力的速度与位置关系}
水平运动受线性阻力 $-kv$ 时，时间方程为 $m\dot v=-kv$。在 $v>0$ 时分离变量，以 $v(0)=v_0>0$ 确定常数：
\[
\int_{v_0}^{v}\frac{\dd w}{w}=-\frac km\int_0^t\dd t',
\qquad
[\ln w]_{v_0}^{v}=-\frac km[t']_0^t,
\qquad
\ln\frac{v}{v_0}=-\frac{kt}{m}.
\]
指数化得到 $v(t)=v_0\e^{-t/\tau}$，$\tau=m/k$。令初始位置为零，再积分速度，
\[
x(t)=v_0\int_0^t\e^{-t'/\tau}\dd t'
=v_0[-\tau\e^{-t'/\tau}]_0^t
=v_0\tau(1-\e^{-t/\tau}).
\]
在同一个 $v>0$ 的运动段，也可写 $mv\dd v/\dd x=-kv$，约去 $v$ 后积分：
$\int_{v_0}^v\dd w=-(k/m)\int_0^x\dd\xi$，
得到 $v-v_0=-kx/m$，即速度随位置线性下降。两种关系描述同一运动，位置线性并不意味着时间线性。

\fig{n4t-drag}{速度随时间指数衰减，随位置则为直线}{\begin{tikzpicture}
\begin{scope}\draw[axis](0,0)--(2.8,0)node[right]{$t$};\draw[axis](0,0)--(0,1.9)node[above]{$v$};\draw[world,domain=0:2.5,samples=40]plot(\x,{1.6*exp(-1.2*\x)});\end{scope}
\begin{scope}[shift={(4,0)}]\draw[axis](0,0)--(2.8,0)node[right]{$x$};\draw[axis](0,0)--(0,1.9)node[above]{$v$};\draw[world](0,1.6)--(2.3,0);\node[below]at(2.3,0){$mv_0/k$};\end{scope}
\end{tikzpicture}}
图\ref{n4t-drag}的直线终点是位置的极限，不是在某个有限时刻突然停住的事件。约去 $v$ 的步骤只在 $v\ne0$ 区间有效；不能将直线继续延长到负速度并解释成阻力把物体反向拉走。模型预测的是无限趋近零速和有限极限位移。

判断是否有限时间到达某位置，可用 $\Delta t=\int\dd x/v(x)$。若接近终点时 $v(x)$ 近似为 $C(x_*-x)^p$，$C>0$，考虑从 $x_*-\varepsilon$ 到 $x_*-\delta$ 的一小段，$0<\delta<\varepsilon$。换元 $y=x_*-x$，有 $\dd x=-\dd y$，上下限从 $\varepsilon$ 变到 $\delta$，故
\[
\Delta t\simeq\frac1C\int_\delta^\varepsilon y^{-p}\dd y
=\begin{cases}
\dfrac{\varepsilon^{1-p}-\delta^{1-p}}{C(1-p)},&p\ne1,\\[3pt]
\dfrac1C\ln(\varepsilon/\delta),&p=1.
\end{cases}
\]
第一行的原函数为 $y^{1-p}/(1-p)$，第二行为 $\ln y$。令 $\delta\to0^+$：$p<1$ 时剩余时间有限，$p\ge1$ 时发散。恒定反向合力给 $v\propto\sqrt{x_*-x}$，所以有限时间停下；线性速度阻力给 $v\propto x_*-x$，则只能无限接近。这里比较的是不同物理阻力模型的终点行为，不能只看空间表达式都存在一个零点。

\section{相对运动、碰撞与重复过程}
\subsection{最近距离由相对速度为零确定}
两物体沿一直线相向或追及，间距变化率是它们的速度差。最近时两者速度相同，但共同速度一般不为零；只有在特定质心系中才可能同时静止。已知相互斥力时，可以分别求绝对加速度，也可以直接写相对加速度 $F(1/m_1+1/m_2)$，再用相对初速求接近距离。

把相互作用写成保守斥力模型时，近距离内的势能储存使相对速度降低，随后又释放，使两物体分开。最近时并不要求外界能量消失。质心保持原来的匀速运动与相对运动的反向，是两个互相独立的描述，应同时满足总动量和位置连续的检查。

\fig{fv:n4:gap}{速度相等给出间距的驻点，最小值仍需与阶段端点比较}{\begin{tikzpicture}[x=1.1cm,y=.9cm]
\draw[->](0,0)--(5.8,0)node[right]{$t$};\draw[->](0,0)--(0,3.5)node[above]{$D$};
\draw[blue,thick,domain=0:4,samples=60]plot(\x,{.7+.35*(\x-2.4)^2});\draw[blue,thick](4,1.596)--(5.1,2.7);
\draw[dashed](2.4,0)--(2.4,.7);\draw[dashed](4,0)--(4,2.3);
\node[below]at(2.4,0){$t_*$};\node[below]at(4,0){$t_1$};\node[right]at(2.4,.7){$\dot D=0$};\end{tikzpicture}}
图\ref{fv:n4:gap}用净间距 $D=x_{\text{前}}-x_{\text{后}}$表示碰撞条件。曲线的斜率才是相对速度；斜率从负转正的 $t_*$为本阶段最小间距点。若算出的 $t_*$已越过停车或碰撞引起的阶段边界 $t_1$，便不是该段可用的驻点，应先比较区间端点，再进入下一段。全程安全要求每一段的最小 $D$都满足题设的间隙条件，不能只检查最初的一次速度相等。


\subsection{碰后交换速度并非全程交换轨迹}
等质量一维完全弹性碰撞在外力冲量可忽略时交换瞬时速度。碰撞后两物体仍各自受到自己的重力分量和摩擦，若摩擦因数不同，后续加速度便不同。因此不能把第一次交换速度后直接理解成两物体从此互换了全部轨迹。

\fig{n4deep:collision}{碰撞只交换边界速度，碰后的斜率仍由各自外力决定}{\begin{tikzpicture}
\draw[axis](0,0)--(5.5,0)node[right]{$t$};\draw[axis](0,0)--(0,3.5)node[above]{$v$};\draw[guide](2,0)--(2,3.3);\node[below]at(2,0){$t_c$};\draw[world,blue](0,1)--(2,3);\draw[world,red](0,1.2)--(2,1.2);\draw[world,blue](2,1.2)--(4,3.2);\draw[world,red](2,3)--(4.7,3);\draw[blue,dashed](2,3)--(2,1.2);\fill[blue](2,1.2)circle(2pt);\fill[red](2,3)circle(2pt);\node[blue]at(4.4,2){$a_A$ 保持};\node[red,above]at(4.5,3){$a_B=0$};\node at(1,-.7){碰前};\node at(3.5,-.7){碰后重新计时};
\end{tikzpicture}}
图\ref{n4deep:collision}以 $A$ 受恒定下滑合力、$B$ 滑动时合力为零为例。碰撞瞬间蓝、红两条线交换高度，表示交换速度；碰后蓝线仍有原来的斜率，红线仍水平。下一次接触由碰后两条速度线之间的有符号面积累计到零确定，不能在图上把两种颜色的整条后续轨迹互换。若有一条线先到零速度，须先处理停车事件。

重复碰撞常由“碰前速度—碰后初值—下一事件”形成递推。先判断下一次是再次接触，还是其中一个物体先停下。若先停下，应将余下过程切成两段；只有事件顺序确定后，才能讨论时间间隔或位移构成等差、等比数列。无穷层、无穷次问题则应先求有限层或有限次结果，再明确所取极限。

速度交换还含质量与弹性两个条件。对两个相同质量的物体，若碰撞期间外力冲量可忽略，内部相互作用按第三定律抵消，故 $v_1^-+v_2^-=v_1^++v_2^+$。完全弹性的一维碰撞还要求相对分离速度等于碰前相对接近速度，即 $v_2^+-v_1^+=v_1^--v_2^-$。两式相加、相减才得到交换速度。若恢复系数小于一，第二式右端还须乘该系数，交换规律就不成立；若两质量不等，也不能用等质量形式的动量关系。

“碰撞很短”本身不是忽略外力冲量的充分表述。还应比较外力乘接触时间同碰撞引起的典型动量变化；当物体速度差已经极小，原来可忽略的重力或摩擦冲量可能变得重要。无穷次碰撞的递推极限正是在这种尺度逐渐缩小时最需谨慎解释。有限次结果仍可严格依照给定理想模型计算，而不必把模型的极限结论扩大成接触材料的真实行为。

\subsection{优化所需时间时保留几何边界}
沿固定光滑直槽下滑，加速度是重力沿槽的分量，行程由槽同目标边界的交点决定。改变槽方向会同时改变加速度与路程，不能只选择最大重力分量。可用几何关系把时间平方写成角度函数，再在允许到达目标的角度区间内取最小值；最优角还应保证槽向下、与目标有正的交会距离。

\fig{n4deep:path}{槽越竖直不一定越快，还须比较它与目标的交点距离}{\begin{tikzpicture}
\coordinate(A)at(0,3);\coordinate(P)at(1.2,.6);\coordinate(B)at(.68,.34);\draw[world](-1,-.5)--(3.4,1.7);\draw[world](A)--(B)node[midway,left]{$s$};\draw[guide](A)--(P)node[midway,right]{$h$};\draw[guide](A)--(0,-.2);\draw[flow,red](-.45,2.4)--(-.45,1.4)node[left]{$\bm g$};\fill(A)circle(2pt)node[above]{$A$};\fill(B)circle(2pt)node[below]{$B$};\node[right]at(P){$P$};\draw(1.11,.78)--(1.29,.87)--(1.38,.69);\node at(.15,2.05){$\beta$};\node at(.56,1.75){$\gamma$};
\end{tikzpicture}}
图\ref{n4deep:path}中 $AP=h$ 为到目标直线的垂距，槽 $AB=s$ 与竖直夹角为 $\beta$，与向下法线 $AP$ 夹角为 $\gamma$。垂直投影给 $s\cos\gamma=h$，受力给 $a_s=g\cos\beta$，于是
\[
t^2=\frac{2s}{a_s}=\frac{2h}{g\cos\beta\cos\gamma}.
\]
优化须同时保留这两个余弦因子，并要求它们都为正；前者保证沿槽向下加速，后者保证槽能正向到达目标。目标倾角固定时，$\beta$ 与 $\gamma$ 还有一个几何关系，代入后才得到单变量极值问题。

这类优化同运动学中的最短时间问题相衔接。几何关系确定路程，受力关系确定沿槽加速度；将它们同时代入时间函数后，应展开所用的三角变换或求导步骤，并把内部极值与允许角度的端点比较，才能确定可实现的最短时间。

\section{例题与分析}
\begin{question}\begin{example}[保持速率的航向改变]\label{n4e1}
质量近似恒为 $m$ 的探测器初速度为 $v_0\bm e_x$，欲改为与正 $x$ 轴向下成 $60^\circ$ 的速度，速率仍为 $v_0$。四发动机 $P_1,P_2,P_3,P_4$ 的推力方向分别为 $-x,+y,+x,-y$，大小均为 $f$，不引起转动。求依次开动的发动机及时间。
\fig{n4e1q}{发动机编号按其所产生的推力方向定义}{\begin{tikzpicture}
\draw[world](-1.4,-.5)rectangle(1.4,.5);\node[right]at(1.4,0){$P_1$};\node[left]at(-1.4,0){$P_3$};\node[above]at(0,.5){$P_4$};\node[below]at(0,-.5){$P_2$};\draw[axis](2.7,0)--(4.2,0)node[right]{$x$};\draw[axis](2.7,0)--(2.7,1.4)node[above]{$y$};\draw[flow](2.7,0)--(3.4,-1.21)node[right]{$\bm v_f$};\node at(3.1,-.3){$60^\circ$};
\end{tikzpicture}}
\end{example}\end{question}
\begin{solution}\textbf{把目标速度分解为所需的两分量增量}\par\nobreak
图\ref{n4e1q}目标速度分量为 $v_{f,x}=v_0/2$、$v_{f,y}=-\sqrt3v_0/2$，所以需要的速度增量是 $\Delta v_x=-v_0/2$、$\Delta v_y=-\sqrt3v_0/2$。应减小正 $x$ 速度，同时增加负 $y$ 速度，因此选 $P_1$ 与 $P_4$。
\fig{n4e1s}{所需推力方向由速度增量而非末速度决定}{\begin{tikzpicture}
\coordinate(O)at(0,0);\coordinate(I)at(2.4,0);\coordinate(F)at(1.2,-2.078);
\draw[flow,blue](O)--(I)node[above]{$\bm v_i$};\draw[flow,red](O)--(F)node[left]{$\bm v_f$};\draw[flow](I)--(F)node[midway,right]{$\Delta\bm v$};\draw[guide](I)--(1.2,0)--(F);
\end{tikzpicture}}


\par\medskip\Needspace{4\baselineskip}\noindent\textbf{逐方向求发动机工作时间}\par\nobreak
由图\ref{n4e1s}两分量的冲量，$ft_1=mv_0/2$、$ft_4=\sqrt3mv_0/2$，所以 $t_1=mv_0/(2f)$、$t_4=\sqrt3mv_0/(2f)$。质量恒定、推力沿固定坐标轴且无其他外力时，两次作用顺序可以互换；只开负 $y$ 发动机不能同时减小 $x$ 分量，末速率反而增大。
\end{solution}

\begin{question}\begin{example}[采样速度与换面时刻]\label{n4e2}
物块由光滑斜面顶端静止下滑，经 $B$ 后进入粗糙水平面，过 $B$ 时速率连续。测得下表数据，取 $g=10\,\mathrm{m/s^2}$，求斜面倾角、水平面摩擦因数及 $t=0.6\,\mathrm s$ 的速度。
\begin{center}\begin{tabular}{c|rrrrr}\toprule
$t/\mathrm s$&0&0.2&0.4&1.2&1.4\\
$v/(\mathrm{m/s})$&0&1.0&2.0&1.1&0.7\\\bottomrule
\end{tabular}\end{center}
\fig{n4e2q}{斜面到水平面的交界不一定对应采样时刻}{\begin{tikzpicture}
\draw[world](-2.5,2)--(0,0)--(3,0);\node[above]at(-2.5,2){$A$};\node[below]at(0,0){$B$};\node[below]at(2.6,0){$C$};\node at(-.6,.2){$\alpha$};\node[draw,rotate=-38.6598,anchor=south]at(-1.7,1.36){物块};
\end{tikzpicture}}
\end{example}\end{question}
\begin{solution}\textbf{从同一运动阶段的采样点求加速度}\par\nobreak
图\ref{n4e2q}前段的三个采样点给 $a_1=(2-1)/(0.4-0.2)=5\,\mathrm{m/s^2}$。光滑斜面有 $g\sin\alpha=a_1$，所以 $\alpha=30^\circ$。后段 $a_2=(0.7-1.1)/(1.4-1.2)=-2\,\mathrm{m/s^2}$，故 $\mu=|a_2|/g=0.2$。
\fig{n4e2s}{两段速度直线在真实换面时刻相交}{\begin{tikzpicture}[x=2.8cm,y=.65cm]
\draw[axis](0,0)--(1.7,0)node[right]{$t/\mathrm s$};\draw[axis](0,0)--(0,3.1)node[above]{$v/(\mathrm{m/s})$};\draw[world](0,0)--(.5,2.5)--(1.5,.5);\draw[guide](.5,0)--(.5,2.5);\node[below]at(.5,0){$0.5$};\fill(.6,2.3)circle(1.4pt)node[above right]{$(0.6,2.3)$};
\end{tikzpicture}}


\par\medskip\Needspace{4\baselineskip}\noindent\textbf{用两速度函数的交点找真实换面时刻}\par\nobreak
斜面速度为 $v=5t$，水平面速度从采样点反推为 $v=1.1-2(t-1.2)=3.5-2t$。图\ref{n4e2s}交点满足 $5t_B=3.5-2t_B$，所以 $t_B=0.5\,\mathrm s$。所求 $0.6\,\mathrm s$ 已进入水平段，$v=3.5-1.2=2.3\,\mathrm{m/s}$。若直接延长斜面公式得到 $3\,\mathrm{m/s}$，就越过了模型切换边界。
\end{solution}

\begin{question}\begin{example}[拉力改变后的速度三角形]\label{n4e3}
质量 $4\,\mathrm{kg}$ 的物体在水平拉力 $F$ 下从静止运动，$10\,\mathrm s$ 时速度达 $10\,\mathrm{m/s}$，此后拉力减为 $F/3$，到 $30\,\mathrm s$ 恰停下。两段速度图线均为直线，取 $g=10\,\mathrm{m/s^2}$，求 $F$ 与动摩擦因数。
\fig{n4e3q}{速度图像给出改变拉力前后的加速度}{\begin{tikzpicture}[x=.12cm,y=.18cm]
\draw[axis](0,0)--(36,0)node[right]{$t/\mathrm s$};\draw[axis](0,0)--(0,13)node[above]{$v/(\mathrm{m/s})$};\draw[world](0,0)--(10,10)--(30,0);\draw[guide](10,0)--(10,10)--(0,10);\node[left]at(0,10){$10$};\node[below]at(10,0){$10$};\node[below]at(30,0){$30$};
\end{tikzpicture}}
\end{example}\end{question}
\begin{solution}\textbf{读出两段斜率并保持摩擦方向}\par\nobreak
图\ref{n4e3q}第一段 $a_1=1\,\mathrm{m/s^2}$，第二段 $a_2=-0.5\,\mathrm{m/s^2}$。两段速度都向右，摩擦均向左且大小相同，记为 $f$，故 $F-f=4\,\mathrm N$、$F/3-f=-2\,\mathrm N$。
\fig{n4e3s}{两段摩擦大小相同，减小的是仍然向右的拉力}{\begin{tikzpicture}
\node[draw,minimum width=.7cm](A)at(0,0){$m$};\draw[force](A.east)--(2.15,0)node[midway,above]{$F$};\draw[force,red](A.west)--(-1.35,0)node[midway,above]{$f$};\draw[flow](-.3,.8)--(.8,.8)node[right]{$v>0$};\node[draw,minimum width=.7cm](B)at(5,0){$m$};\draw[force](B.east)--(5.95,0)node[above right]{$F/3$};\draw[force,red](B.west)--(3.65,0)node[midway,above]{$f$};\draw[flow](4.7,.8)--(5.8,.8)node[right]{$v>0$};
\end{tikzpicture}}


\par\medskip\Needspace{4\baselineskip}\noindent\textbf{消去摩擦求拉力，再求摩擦因数}\par\nobreak
图\ref{n4e3s}对应的两式相减，$2F/3=6\,\mathrm N$，所以 $F=9\,\mathrm N$、$f=5\,\mathrm N$，$\mu=f/(mg)=5/40=0.125$。停下以后还须用静摩擦条件判断是否保持静止，不能把第二段直线无限延长为反向运动。
\end{solution}

\begin{question}\begin{example}[力随位置变化与等位移分段]\label{n4e4}
质量 $m$ 的物体初速 $v_0\ge0$，沿正 $x$ 方向运动，水平合力从 $x=0$ 处的 $F_0$ 线性变为 $x=l$ 处的 $F_l$。求末速。再讨论：把总位移等分成 $n$ 段，每段合力恒定，各段力大小组成首项 $F_0$、末项 $F_l$ 的等差数列时，末速是否相同。
\fig{n4e4q}{连续直线与等位移阶梯力采用不同的局部模型}{\begin{tikzpicture}
\draw[axis](0,0)--(4.8,0)node[right]{$x$};\draw[axis](0,0)--(0,2.7)node[above]{$F$};\draw[world](0,.7)--(4,2.1);\draw[guide](4,0)--(4,2.1);\node[left]at(0,.7){$F_0$};\node[right]at(4,2.1){$F_l$};\node[below]at(4,0){$l$};
\end{tikzpicture}}
\end{example}\end{question}
\begin{solution}\textbf{把连续合力曲线下的面积转换为动能变化}\par\nobreak
图\ref{n4e4q}描述的是合力随位移的函数，应该累积功。直线力函数为
$F(x)=F_0+(F_l-F_0)x/l$，由 $mv\dd v=F(x)\dd x$ 得
\begin{align*}
m\int_{v_0}^{v_l}w\dd w
&=\int_0^l\left[F_0+\frac{F_l-F_0}{l}x\right]\dd x,\\
\frac m2(v_l^2-v_0^2)
&=\left[F_0x+\frac{F_l-F_0}{2l}x^2\right]_0^l\\
&=F_0l+\frac{F_l-F_0}{2}l=\frac{F_0+F_l}{2}l.
\end{align*}
这与图中的梯形面积相符，因此
\[v_l=\sqrt{v_0^2+\frac{(F_0+F_l)l}{m}}.\]
\fig{n4e4s}{等位移分段的功等于每段矩形面积之和}{\begin{tikzpicture}
\draw[axis](0,0)--(5,0)node[right]{$x$};\draw[axis](0,0)--(0,2.5)node[above]{$F$};\foreach \i/\h in {0/.7,1/1.166667,2/1.633333,3/2.1}{\fill[blue!10](\i,0)rectangle({\i+1},\h);\draw[world](\i,\h)--({\i+1},\h);}\node[below]at(4,0){$l$};\node at(2,2.3){$\Delta x=l/n$};
\end{tikzpicture}}


\par\medskip\Needspace{4\baselineskip}\noindent\textbf{对等位移段求和并检查终点可达性}\par\nobreak
图\ref{n4e4s}分段情形不必取极限。取 $n\ge2$，第 $j$段两端速度记为 $v_{j-1},v_j$，每段精确满足 $v_j^2-v_{j-1}^2=2F_j(l/n)/m$。相加后中间速度平方逐项消去：
\[
 (v_1^2-v_0^2)+\cdots+(v_n^2-v_{n-1}^2)
 =v_n^2-v_0^2=\frac{2l}{mn}\sum_{j=1}^nF_j.
\]
等差序列的第 $j$项为 $F_j=F_0+(j-1)(F_l-F_0)/(n-1)$，因此
\begin{align*}
 \sum_{j=1}^nF_j
 &=nF_0+\frac{F_l-F_0}{n-1}\sum_{j=1}^n(j-1)\\
 &=nF_0+\frac{F_l-F_0}{n-1}\frac{n(n-1)}2
 =\frac{n(F_0+F_l)}2.
\end{align*}
代回后 $v_n^2-v_0^2=(F_0+F_l)l/m$，所以末速仍为同一表达式。这里相同来自等位移权重和等差力序列，不表示两个过程的 $v(t)$或所需时间相同。$n=1$时只有一个力值，首末项必须相同，另按恒力处理。若 $F$是驱动力而非合力，还须另外扣除摩擦功；若根号内虽为正但中途先停下，也要检查该终点是否实际可达。
\end{solution}

\begin{question}\begin{example}[由撤力后的时间反推驱动力]\label{n4e5}
质量 $1.5\,\mathrm{kg}$ 的物块在水平恒力 $F$ 下从 $A$ 静止出发，撤力后再滑行 $2\,\mathrm s$ 停在 $B$。总位移 $5\,\mathrm m$，动摩擦因数 $0.2$，取 $g=10\,\mathrm{m/s^2}$。求 $F$ 及撤力时的位置和速度。
\fig{n4e5q}{撤力事件把全程分成两段}{\begin{tikzpicture}
\draw[world](-3,0)--(3,0);\fill(-2.5,0)circle(2pt)node[below]{$A$};\fill(-1.5,0)circle(2pt)node[below]{$C$};\fill(2.5,0)circle(2pt)node[below]{$B$};\draw[force](-2.5,.6)--(-1.5,.6)node[above]{$F$};\node at(.5,.7){撤力后 $2\,\mathrm s$};\draw[dim](-2.5,-.7)--(2.5,-.7)node[midway,below]{$5\,\mathrm m$};
\end{tikzpicture}}
\end{example}\end{question}
\begin{solution}\textbf{从撤力后的减速段反推速度和位移}\par\nobreak
图\ref{n4e5q}撤力后减速度大小为 $\mu g=2\,\mathrm{m/s^2}$，所以撤力时 $v_C=2\times2=4\,\mathrm{m/s}$，此后位移 $s_2=(v_C+0)\times2/2=4\,\mathrm m$。前段位移因而只有 $s_1=5-4=1\,\mathrm m$。
\fig{n4e5s}{速度图下两块三角形面积对应两段位移}{\begin{tikzpicture}
\draw[axis](0,0)--(5,0)node[right]{$t$};\draw[axis](0,0)--(0,2.3)node[above]{$v$};\fill[blue!10](0,0)--(.8,1.8)--(.8,0)--cycle;\fill[red!7](.8,0)--(.8,1.8)--(4,0)--cycle;\draw[world](0,0)--(.8,1.8)--(4,0);\node at(.45,.35){$s_1$};\node at(2,.45){$s_2$};\node[above]at(.8,1.8){$4\,\mathrm{m/s}$};
\end{tikzpicture}}


\par\medskip\Needspace{4\baselineskip}\noindent\textbf{用剩余位移求前段加速度及驱动力}\par\nobreak
前段由静止匀加速，$v_C^2=2a_1s_1$，故 $a_1=8\,\mathrm{m/s^2}$。图\ref{n4e5s}前段受力方程 $F-\mu mg=ma_1$ 给 $F=1.5(8+2)=15\,\mathrm N$。前段时间 $v_C/a_1=0.5\,\mathrm s$，图像面积 $4\times0.5/2=1\,\mathrm m$，与位置计算一致。
\end{solution}

\begin{question}\begin{example}[飞行器失去与恢复升力]\label{n4e6}
飞行器质量 $2\,\mathrm{kg}$，动力提供向上恒力 $28\,\mathrm N$，运动阻力大小恒定且反向于速度。第一次由地面静止竖直上升，$8\,\mathrm s$ 升高 $64\,\mathrm m$。第二次在上升 $6\,\mathrm s$ 后失去升力。取 $g=10\,\mathrm{m/s^2}$，求阻力、最大高度，以及从开始下落到恢复升力的最长允许延迟。
\fig{n4e6q}{上升、无动力下降、恢复动力制动的受力不同}{\begin{tikzpicture}
\draw[axis](0,0)--(0,3.7)node[above]{$z$};\node[draw]at(.8,1.2){$m=2\,\mathrm{kg}$};\draw[force,blue](.8,1.5)--(.8,2.8)node[above]{$28\,\mathrm N$};\draw[guide](0,3.1)--(3.3,3.1)node[right]{$h$};\draw[world](-.5,0)--(3.5,0);
\end{tikzpicture}}
\end{example}\end{question}
\begin{solution}\textbf{用第一次上升数据确定阻力}\par\nobreak
按图\ref{n4e6q}的竖直过程，第一次的上升加速度为 $a_1=2H/t^2=2\,\mathrm{m/s^2}$，受力 $28-20-f=2\times2$，故 $f=4\,\mathrm N$。

\par\medskip\Needspace{4\baselineskip}\noindent\textbf{计算失去升力后还能上升的高度}\par\nobreak
第二次失去动力时高度 $s_1=2\times6^2/2=36\,\mathrm m$，向上速度 $v_1=12\,\mathrm{m/s}$。继续上升时重力与阻力都向下，减速度 $a_2=(20+4)/2=12\,\mathrm{m/s^2}$，再上升 $s_2=12^2/(2\times12)=6\,\mathrm m$，所以最高点 $h=42\,\mathrm m$。
\fig{n4e6s}{下降时阻力向上，恢复升力后与动力共同制动}{\begin{tikzpicture}
\foreach \x/\s in {-1.7/无动力下降,1.7/恢复动力}{\fill(\x,0)circle(2pt);\draw[force,red](\x,0)--(\x,-1.3)node[below]{$mg$};\draw[force,blue](\x,0)--(\x,1)node[above]{$f$};\node at(\x,-1.9){\s};}\draw[force](2.1,0)--(2.1,1.8)node[above]{$F$};
\end{tikzpicture}}



\par\medskip\Needspace{4\baselineskip}\noindent\textbf{按下降方向重画受力并求恢复时状态}\par\nobreak
自最高点下落后，图\ref{n4e6s}左侧向下加速度 $a_3=(20-4)/2=8\,\mathrm{m/s^2}$。延迟 $t$ 恢复动力，已下降 $4t^2$、向下速率 $8t$。恢复时还在下降，阻力也向上，因此制动加速度为 $a_4=(28+4-20)/2=6\,\mathrm{m/s^2}$，不是上升时的 $2$。



\par\medskip\Needspace{4\baselineskip}\noindent\textbf{把延迟下落与制动下落合并限制}\par\nobreak
\fig{n4deep:brakingheight}{恢复升力时剩余高度还必须容纳随后制动距离}{\begin{tikzpicture}
\draw[axis](0,0)--(0,4.7)node[above]{$z$};\draw[world](-.5,0)--(4.2,0);\draw[guide](0,4)--(3.4,4);\node[left]at(0,4){$42\,\mathrm m$};\fill(1,4)circle(2pt)node[above]{最高点};\fill(1,2.7)circle(2pt)node[right]{恢复升力};\fill(1,.7)circle(2pt)node[right]{制动结束};\draw[flow,blue](1,4)--(1,2.7);\draw[flow,red](1,2.7)--(1,.7);\draw[dim](3,4)--(3,2.7)node[midway,right]{$4t^2$};\draw[dim](3,2.7)--(3,.7)node[midway,right]{$16t^2/3$};\node at(1.5,-.6){制动结束位置须高于地面};
\end{tikzpicture}}
图\ref{n4deep:brakingheight}中，恢复动力是中间事件，物体此时仍有向下速度；允许延迟由两段下降距离之和确定。

制动到零还需向下距离 $(8t)^2/(2\times6)=16t^2/3$。不触地要求
\[
 4t^2+\frac{16t^2}{3}<42
 \quad\Longleftrightarrow\quad \frac{28}{3}t^2<42
 \quad\Longleftrightarrow\quad t^2<\frac92.
\]
由于 $t\ge0$，得到 $t<3/\sqrt2\,\mathrm s$。临界延迟约为 $2.12\,\mathrm s$；严格不触地时它是上确界，不能在等号处取得，因为等号使最低点恰在地面。恢复后停住若仍有动力，下一段重新上升，加速度才变回 $2\,\mathrm{m/s^2}$。
\end{solution}

\begin{question}\begin{example}[不同摩擦下的重复弹性碰撞]\label{n4e7}
倾角 $30^\circ$ 的长斜面上，两等质量物块 $A,B$ 初静止，$A$ 在上、间距 $L=0.3\,\mathrm m$。取 $g=10\,\mathrm{m/s^2}$，$\mu_A=2\sqrt3/9$，各处最大静摩擦等于动摩擦。每次碰撞时间忽略、两块交换速度。求第一次碰后 $B$ 速率，并分别对 $\mu_B=\sqrt3/3$、$\sqrt3/2$ 求第 $n$ 次碰撞时 $A$ 的累计路程。
\fig{n4e7q}{碰撞只交换速度，之后仍各受自己的摩擦}{\begin{tikzpicture}
\draw[world](-2.5,0)--(2.5,{5*tan(30)});\node[draw,rotate=30,anchor=south]at(1.3,{3.8*tan(30)}){$A$};\node[draw,rotate=30,anchor=south]at(-.5,{2*tan(30)}){$B$};\draw[dim](-.2,{2*tan(30)-.52})--(1.6,{3.8*tan(30)-.52})node[midway,below]{$L$};\node at(-1.6,.2){$30^\circ$};
\end{tikzpicture}}
\end{example}\end{question}
\begin{solution}\textbf{求第一次碰前速度并应用速度交换}\par\nobreak
沿图\ref{n4e7q}斜面向下取正。$A$ 的滑动加速度 $a=g(\sin30^\circ-\mu_A\cos30^\circ)=5/3\,\mathrm{m/s^2}$。两种 $\mu_B$ 都足以让 $B$ 初始静止，所以第一次接触前 $A$ 速度 $u_1=\sqrt{2aL}=1\,\mathrm{m/s}$，碰后 $B$ 得到此速度，$A$ 瞬时速度为零。
\fig{n4e7s}{第一次碰后的速度线明确显示追撞与停车的先后}{\begin{tikzpicture}
\begin{scope}[x=2.3cm,y=1cm]\draw[axis](0,0)--(1.4,0)node[right]{$\tau/\mathrm s$};\draw[axis](0,0)--(0,2.6)node[above]{$v$};\draw[world,blue](0,0)--(1.2,2)node[right]{$A$};\draw[world,red](0,1)--(1.2,1)node[right]{$B$};\draw[guide](1.2,0)--(1.2,2);\node[below]at(1.2,0){$1.2$};\node[left]at(0,1){$1$};\node at(.6,3){$B$ 匀速};\end{scope}
\begin{scope}[xshift=5.5cm,x=5.5cm,y=2cm]\draw[axis](0,0)--(.62,0)node[right]{$\tau/\mathrm s$};\draw[axis](0,0)--(0,1.3)node[above]{$v$};\draw[world,blue](0,0)--(.489898,.816497)node[above]{$A$};\draw[world,red](0,1)--(.4,0)--(.55,0);\draw[guide](.489898,0)--(.489898,.816497);\node[below left]at(.4,0){$.4$};\draw[guide](.489898,0)--(.52,-.25)node[below]{$\sqrt{.24}$};\node[left]at(0,1){$1$};\node at(.28,1.5){$B$ 先停后被追上};\end{scope}
\end{tikzpicture}}



\par\medskip\Needspace{4\baselineskip}\noindent\textbf{不减速情形用相对位移求碰撞间隔}\par\nobreak
先取 $\mu_B=\sqrt3/3$，$B$ 一旦滑动便匀速。第一次碰后 $A$ 从零加速，$B$ 以 $u_1=1$ 匀速，下一次接触满足 $a\tau^2/2=u_1\tau$，故间隔 $\tau=2u_1/a=1.2\,\mathrm s$。第二次碰前速度为 $2,1$，交换后成为 $1,2$。此后每次速度差仍为 $1\,\mathrm{m/s}$，所以间隔恒为 $1.2\,\mathrm s$，第 $j$ 次碰后 $B$ 速度为 $j\,\mathrm{m/s}$。到第 $n$ 次碰撞，$B$ 位移为 $1.2\sum_{j=1}^{n-1}j=0.6n(n-1)\,\mathrm m$。两块碰撞时在同一位置，因此
\[s_{A,n}=0.3+0.6n(n-1)\quad(\mathrm m).\]



\par\medskip\Needspace{4\baselineskip}\noindent\textbf{强摩擦情形先验证停车先于追上}\par\nobreak
再取 $\mu_B=\sqrt3/2$，运动中的 $B$ 以 $b=2.5\,\mathrm{m/s^2}$ 减速。设某次碰后 $A$ 静止、$B$ 速度 $u_j$。$B$ 停下需 $u_j/b$，移动 $d_j=u_j^2/(2b)$；在这段时间 $A$ 只移动 $a u_j^2/(2b^2)=(a/b)d_j<d_j$。所以图\ref{n4e7s}右图的事件顺序确为 $B$ 先停下，随后 $A$ 再追到它。



\par\medskip\Needspace{4\baselineskip}\noindent\textbf{建立速度递推并累加碰撞位置}\par\nobreak
下次碰前 $A$走过 $d_j$，速度 $u_{j+1}=\sqrt{2ad_j}=u_j\sqrt{a/b}=u_j\sqrt{2/3}$；$B$此时静止，交换速度后又回到同类初态。将递推式平方，$u_{j+1}^2=(2/3)u_j^2$，由 $u_1=1\,\mathrm{m/s}$得到 $u_j^2=(2/3)^{j-1}\,\mathrm{m^2/s^2}$。再用 $d_j=u_j^2/(2b)$，有 $d_j=0.2(2/3)^{j-1}\,\mathrm m$。从首次接触位置开始累计：
\begin{align*}
 s_{A,n}&=L+\sum_{j=1}^{n-1}d_j\\
 &=0.3+0.2\frac{1-(2/3)^{n-1}}{1-2/3}
 \quad(\mathrm m).
\end{align*}
等比求和中的分母为 $1/3$，所以
\[s_{A,n}=0.3+0.6\left[1-(2/3)^{n-1}\right]\quad(\mathrm m).\]


\par\medskip\Needspace{4\baselineskip}\noindent\textbf{检查首项和有限路程极限}\par\nobreak
两式在 $n=1$ 都回到 $0.3\,\mathrm m$。后一模型的无穷次碰撞路程趋于 $0.9\,\mathrm m$；它是理想瞬时碰撞与摩擦模型的极限，不把无限延续的微小碰撞直接解释为真实接触细节。
\end{solution}

\begin{question}\begin{example}[分段拉力下长板与物块的速度追平]\label{n4e8}
质量 $m$ 的物块放在质量 $2m$ 的长板上，地面光滑，层间动摩擦因数 $0.2$，最大静摩擦近似等于动摩擦。整体初静止，拉板的右向力在 $0$--$1\,\mathrm s$ 为 $mg$，$1$--$1.5\,\mathrm s$ 为 $0.4mg$，之后为零。取 $g=10\,\mathrm{m/s^2}$，求关键时刻两速度及 $0$--$3\,\mathrm s$ 的相对滑动距离。
\fig{n4e8q}{拉力撤去后，板与块仍可继续相对滑动}{\begin{tikzpicture}
\node[draw,minimum width=3cm,minimum height=.6cm]at(-1,.3){长板 $2m$};\node[draw,anchor=south]at(-.2,.6){$m$};\draw[force](.5,.3)--(1.7,.3)node[right]{$F$};
\begin{scope}[shift={(3,0)}]\draw[axis](0,0)--(3.1,0)node[right]{$t$};\draw[axis](0,0)--(0,1.9)node[above]{$F/(mg)$};\draw[world](0,1.6)--(1,1.6);\draw[world](1,.64)--(1.5,.64);\draw[world](1.5,0)--(2.8,0);\draw[guide](1,0)--(1,1.6) (1.5,0)--(1.5,.64);\node[below]at(1,0){$1$};\node[below]at(1.5,0){$1.5$};\end{scope}
\end{tikzpicture}}
\end{example}\end{question}
\begin{solution}\textbf{检验起始阶段能否共同运动}\par\nobreak
共同运动所能提供的最大驱动力为 $(3m)\mu g=0.6mg$，所以图\ref{n4e8q}第一段必滑动。上块向右加速度 $a_m=\mu g=2$，长板 $a_P=(mg-0.2mg)/(2m)=4$，单位为 $\mathrm{m/s^2}$。$t=1$ 时 $v_m=2$、$v_P=4\,\mathrm{m/s}$。



\par\medskip\Needspace{4\baselineskip}\noindent\textbf{沿既有滑动方向逐段更新速度}\par\nobreak
第二段虽然外力降低，但板仍比块快，摩擦方向不变。此段 $a_m=2$、$a_P=(0.4mg-0.2mg)/(2m)=1$，到 $1.5\,\mathrm s$ 有 $v_m=3$、$v_P=4.5\,\mathrm{m/s}$。之后外力为零，仍在滑动时 $a_m=2$、$a_P=-1$，速度差以 $3\,\mathrm{m/s^2}$ 减小，从 $1.5$ 降到零需 $0.5\,\mathrm s$。
\fig{n4e8s}{两速度图之间的面积是相对位移}{\begin{tikzpicture}[x=1.35cm,y=.46cm]
\draw[axis](0,0)--(3.5,0)node[right]{$t/\mathrm s$};\draw[axis](0,0)--(0,5.7)node[above]{$v/(\mathrm{m/s})$};\fill[blue!10](0,0)--(1,4)--(1.5,4.5)--(2,4)--cycle;\draw[world,blue](0,0)--(1,4)--(1.5,4.5)--(2,4)--(3,4);\draw[world,red](0,0)--(2,4)--(3,4);\node[above]at(1.25,4.6){板};\node[right]at(.8,1.6){块};\foreach \t in {1,1.5,2,3}{\draw[guide](\t,0)--(\t,4);\node[below]at(\t,0){$\t$};}
\end{tikzpicture}}



\par\medskip\Needspace{4\baselineskip}\noindent\textbf{速度追平后检验共同运动能否维持}\par\nobreak
在 $t=2\,\mathrm s$ 两者均为 $4\,\mathrm{m/s}$。此时若共同运动，所需摩擦为零，静摩擦可以维持，故 $t=3\,\mathrm s$ 两速度仍为 $4\,\mathrm{m/s}$。长板在 $1,1.5,2,3\,\mathrm s$ 的速度依次为 $4,4.5,4,4\,\mathrm{m/s}$。



\par\medskip\Needspace{4\baselineskip}\noindent\textbf{用速度差面积计算实际滑动距离}\par\nobreak
图\ref{n4e8s}中直到 $2\,\mathrm s$板始终比块快，故相对滑动距离等于速度差的积分。以下时间以秒、速度以米每秒计；每段另取当地时间 $\tau$，保留该段的初始速度差：
\begin{align*}
 s_1&=\int_0^1 2\tau\dd\tau=[\tau^2]_0^1=1\,\mathrm m,\\
 s_2&=\int_0^{0.5}(2-\tau)\dd\tau
 =\left[2\tau-\frac{\tau^2}{2}\right]_0^{0.5}
 =0.875\,\mathrm m,\\
 s_3&=\int_0^{0.5}(1.5-3\tau)\dd\tau
 =\left[1.5\tau-\frac32\tau^2\right]_0^{0.5}
 =0.375\,\mathrm m.
\end{align*}
共同运动后的速度差为零，所以 $0$--$3\,\mathrm s$总滑动距离为 $s_1+s_2+s_3=2.25\,\mathrm m$。达到速度相同时才切换到共同运动，不能在撤力的 $1.5\,\mathrm s$就提前停止计算滑动。
\end{solution}

\begin{question}\begin{example}[两变力驱动相邻物块的分离]\label{n4e10}
光滑水平面上左块 $A$ 质量 $1\,\mathrm{kg}$，右块 $B$ 质量 $2\,\mathrm{kg}$，初静止并接触。两者各受右向外力 $F_1=(9-2t)\,\mathrm N$、$F_2=(3+2t)\,\mathrm N$，$t$ 以秒计。求分离时刻、此时加速度、分离后两秒内的加速度图像及两秒后的相对速度。
\fig{n4e10q}{两个外力都向右，压力能否保持取决于各自加速需求}{\begin{tikzpicture}
\draw[world](-3,0)--(3,0);\node[draw,minimum width=1.3cm,minimum height=.6cm,anchor=south]at(-.65,0){$A,1$};\node[draw,minimum width=1.3cm,minimum height=.6cm,anchor=south]at(.65,0){$B,2$};\draw[force](-2.7,.3)--(-1.3,.3)node[midway,above]{$F_1$};\draw[force](1.3,.3)--(2.8,.3)node[midway,above]{$F_2$};
\end{tikzpicture}}
\end{example}\end{question}
\begin{solution}\textbf{接触阶段由整体加速度求内部压力}\par\nobreak
图\ref{n4e10q}接触期间整体受力 $F_1+F_2=12\,\mathrm N$，共同加速度 $a=12/3=4\,\mathrm{m/s^2}$。隔离 $A$，$F_1-N=(1\,\mathrm{kg})a$，故 $N=(5-2t)\,\mathrm N$，到 $t_s=2.5\,\mathrm s$ 恰为零。在此之前非负，之后不能再继续共同加速。
\fig{n4e10s}{以分离时刻为局部时间零点画加速度}{\begin{tikzpicture}[x=1.7cm,y=.35cm]
\draw[axis](0,0)--(2.6,0)node[right]{$\tau/\mathrm s$};\draw[axis](0,0)--(0,7.5)node[above]{$a/(\mathrm{m/s^2})$};\fill[blue!10](0,4)--(2,6)--(2,0)--cycle;\draw[world,blue](0,4)--(2,6)node[right]{$a_B$};\draw[world,red](0,4)--(2,0)node[above right]{$a_A$};\node[left]at(0,4){$4$};\node[below]at(2,0){$2$};
\end{tikzpicture}}


\par\medskip\Needspace{4\baselineskip}\noindent\textbf{以分离时刻为零点重建两加速度}\par\nobreak
令 $\tau=t-2.5$，图\ref{n4e10s}分离后 $a_A=4-2\tau$、$a_B=4+\tau$，单位为 $\mathrm{m/s^2}$。两者在 $\tau=0$ 都为 $4$，相对加速度为 $3\tau$，故 $B$ 逐渐领先，符合分离假设。



\par\medskip\Needspace{4\baselineskip}\noindent\textbf{积分相对加速度并检验间隙增长}\par\nobreak
加速度图下的有符号面积为速度增量，两曲线间面积为相对速度增量。分离时相对速度和间隙均为零，因而
\[
v_B-v_A=\int_0^\tau3s\,\dd s=[3s^2/2]_0^\tau
=\frac32\tau^2.
\]
代入 $\tau=2\,\mathrm s$ 得 $6\,\mathrm{m/s}$。再积分一次，
\[
\Delta x=\int_0^\tau\frac32s^2\,\dd s
=[s^3/2]_0^\tau=\frac{\tau^3}{2},
\]
在此段始终增加，不会重新接触。
\end{solution}

\begin{question}\begin{example}[有限作用范围内的恒定斥力]\label{n4e11}
光滑直线上，视为质点的 $A$ 质量 $m_1=1\,\mathrm{kg}$，初静止；质点 $B$ 质量 $m_2=3\,\mathrm{kg}$，以 $v_0=0.20\,\mathrm{m/s}$ 从左侧接近。间距不小于 $d=0.10\,\mathrm m$ 时无相互作用，小于 $d$ 时有大小 $F=0.60\,\mathrm N$ 的相互斥力。求相互作用阶段两加速度、最近时两速度和最小间距。
\fig{n4e11q}{相互作用从间距达到 $d$ 后开始}{\begin{tikzpicture}
\draw[world](-3,-.1)--(3,-.1);\fill(-1.7,0)circle(3pt)node[above]{$B,3$};\fill(1.4,0)circle(3pt)node[above]{$A,1$};\draw[flow](-1.7,1)--(-.4,1)node[right]{$v_0$};\draw[dim](-1.7,-.5)--(1.4,-.5)node[midway,below]{$d$ 时开始计时};
\end{tikzpicture}}
\end{example}\end{question}
\begin{solution}\textbf{进入作用范围后求速度相等时刻}\par\nobreak
在图\ref{n4e11q}作用开始后取向右为正，$a_A=F/m_1=0.60\,\mathrm{m/s^2}$，$a_B=-F/m_2=-0.20\,\mathrm{m/s^2}$。局部时间 $\tau$ 内，$v_A=0.60\tau$、$v_B=0.20-0.20\tau$。最近条件是两速度相等，故 $\tau_*=0.25\,\mathrm s$，共同速度 $0.15\,\mathrm{m/s}$。
\fig{n4e11s}{最近时相对速度为零，质心仍向右运动}{\begin{tikzpicture}
\draw[axis](0,0)--(4.8,0)node[right]{$\tau$};\draw[axis](0,0)--(0,2.5)node[above]{$v$};\draw[world,blue](0,0)--(4,2.4)node[right]{$v_A$};\draw[world,red](0,1.6)--(4,.8)node[right]{$v_B$};\fill(2,1.2)circle(2pt);\draw[guide](2,0)--(2,1.2);\node[below]at(2,0){$0.25\,\mathrm s$};
\end{tikzpicture}}


\par\medskip\Needspace{4\baselineskip}\noindent\textbf{计算接近距离并用总动量核查共同速度}\par\nobreak
按图\ref{n4e11s}，从作用开始到速度相等时，接近距离为速度差的积分：
\[
\int_0^{0.25}(0.20-0.80\tau)\dd\tau
=[0.20\tau-0.40\tau^2]_0^{0.25}
=0.050-0.025=0.025\,\mathrm m.
\]
因此最小间距 $d_{\min}=0.10-0.025=0.075\,\mathrm m$。总动量给共同速度 $3\times0.20/(1+3)=0.15\,\mathrm{m/s}$，与逐段计算一致。最近时两球都未停下，只有相对运动暂时停止；随后斥力使它们重新远离。
\end{solution}

\begin{question}\begin{example}[从定点到斜面的最短下滑时间]\label{n4e12}
定点 $A$ 位于倾角为 $\alpha$ 的直斜面上方，垂直于斜面的距离为 $h$。从 $A$ 沿一条光滑直槽静止下滑到斜面，直槽与竖直向下方向向右夹角为 $\beta$。求最短时间对应的槽方向，$0<\alpha<\pi/2$。
\fig{n4e12q}{改变槽方向会同时改变路程与沿槽加速度}{\begin{tikzpicture}
\draw[world](-2,0)--(3,2);\coordinate(A)at(-1,2.5);\coordinate(B)at(-.624,.55);\coordinate(P)at(-.276,.69);
\fill(A)circle(2pt)node[above]{$A$};\draw[world](A)--(B)node[below]{$B$};\draw[guide](A)--(P)node[midway,right]{$h$};\draw[guide](A)--(-1,.2);\node at(-.87,1.75){$\beta$};\draw[guide](-2,0)--(0,0);\node at(-1.2,.15){$\alpha$};
\end{tikzpicture}}
\end{example}\end{question}
\begin{solution}\textbf{由斜面法向投影求槽长并保留可达条件}\par\nobreak
图\ref{n4e12q}槽方向为 $(\sin\beta,-\cos\beta)$，同指向斜面的法向夹角为 $\alpha-\beta$。因此路径长 $s=h/\cos(\alpha-\beta)$，必须 $\cos(\alpha-\beta)>0$。沿槽受重力分量，$a=g\cos\beta>0$，由 $s=at^2/2$ 得
\[t^2=\frac{2h}{g\cos\beta\cos(\alpha-\beta)}.\]
\fig{n4e12s}{最优槽方向平分竖直向下与斜面向下法线的夹角}{\begin{tikzpicture}
\coordinate(A)at(0,3);\draw[flow](A)--++(-90:3)node[below left]{竖直};\draw[flow](A)--++(-50:3)node[right]{斜面法向};\draw[flow,blue](A)--++(-70:3)node[below right]{$\beta=\alpha/2$};\draw(0,2.1)arc(-90:-70:.9);\draw({1.35*cos(-70)},{3+1.35*sin(-70)})arc(-70:-50:1.35);\node at(.18,1.8){$\beta$};\node at(1.05,1.3){$\alpha-\beta$};
\end{tikzpicture}}


\par\medskip\Needspace{4\baselineskip}\noindent\textbf{把时间最小化转为余弦乘积最大化}\par\nobreak
利用 $\cos\beta\cos(\alpha-\beta)=[\cos\alpha+\cos(2\beta-\alpha)]/2$，最大值在 $2\beta-\alpha=0$ 取得。所以图\ref{n4e12s}最优 $\beta=\alpha/2$，$t_{\min}=\sqrt{2h/g}/\cos(\alpha/2)$。该角位于可达区间内部，槽向下且同斜面正向相交。竖直槽虽有最大加速度，却未同时取得最合适的路径长度。
\end{solution}

\begin{question}\begin{example}[无穷叠层楔块的有限层极限]\label{n4e13}
每一层由两个质量均为 $m$ 的直角三角块拼成正方形，斜接触面与水平成 $45^\circ$，所有斜面方向相同；相邻层间为水平光滑接触，底层在光滑地面上。各块由静止释放。先求 $n$ 层时底部支持力的递推关系，再定义并求无穷层极限。
\fig{n4e13q}{把无穷叠层定义为有限层初始动力学的极限}{\begin{tikzpicture}
\foreach \j in {0,1,2}{\draw[world](0,{1.4*\j})rectangle(1.4,{1.4*(\j+1)});\draw[world](0,{1.4*\j})--(1.4,{1.4*(\j+1)});\node at(.35,{1.4*\j+.95}){$m$};\node at(1,{1.4*\j+.4}){$m$};}\draw[guide](.3,4.2)--(.3,5) (1.1,4.2)--(1.1,5);\draw[world](-.7,-.05)--(2.1,-.05);\draw[force,blue](.7,-1.2)--(.7,-.05)node[right]{$N_n$};\node at(2.6,2.5){共 $n$ 层};
\end{tikzpicture}}
\end{example}\end{question}
\begin{solution}\textbf{用有限层支持力定义递推量}\par\nobreak
定义图\ref{n4e13q}中 $N_n=c_nmg$，空叠层取 $c_0=0$。分析最下面一层：下三角块只能水平运动，上三角块向下加速度为 $a$，它承受上方 $n-1$ 层的压力 $N_1$。上方系统的底部随它向下加速，内部有效重力由 $g$ 改为 $g-a$，所以 $N_1=c_{n-1}m(g-a)$。
\fig{n4e13s}{底层的两个水平加速度相反，竖直下降量由斜面约束给出}{\begin{tikzpicture}
\draw[world](0,0)rectangle(2,2);\draw[world](0,0)--(2,2);\draw[flow](.5,1.5)--(-.8,1.5)node[left]{$a/2$};\draw[flow](1.5,.5)--(2.8,.5)node[right]{$a/2$};\draw[flow](.7,1.6)--(.7,.4)node[below]{$a$};\draw[force,red](1,2.9)--(1,2)node[midway,right]{$N_1$};\draw[force,blue](1,-1)--(1,0)node[midway,right]{$N_n$};
\end{tikzpicture}}



\par\medskip\Needspace{4\baselineskip}\noindent\textbf{由底层几何约束关联三方向受力}\par\nobreak
两块水平只受相互作用，等质量给反向等大的水平加速度；$45^\circ$ 约束使图\ref{n4e13s}这两个大小均为 $a/2$。设斜面法向力为 $T$，则 $T/\sqrt2=ma/2$。上块竖直式为 $mg+N_1-T/\sqrt2=ma$，下块竖直式为 $N_n=mg+T/\sqrt2$。



\par\medskip\Needspace{4\baselineskip}\noindent\textbf{消去加速度建立有限层递推}\par\nobreak
令 $c=c_{n-1}$，代入得 $(1+c)g=(c+3/2)a$，从而
\[c_n=1+\frac{a}{2g}=\frac{3c_{n-1}+4}{2c_{n-1}+3},\qquad c_0=0.\]


\par\medskip\Needspace{4\baselineskip}\noindent\textbf{证明有界收敛再求无限层极限}\par\nobreak
第一层为 $c_1=4/3$。记映射为 $f(c)=(3c+4)/(2c+3)$，在 $0\le c<\sqrt2$时分别检查上界和递增性：
\begin{align*}
 \sqrt2-f(c)
 &=\frac{(2\sqrt2-3)c+3\sqrt2-4}{2c+3}
 =\frac{(3-2\sqrt2)(\sqrt2-c)}{2c+3}>0,\\
 f(c)-c&=\frac{3c+4-c(2c+3)}{2c+3}
 =\frac{4-2c^2}{2c+3}>0.
\end{align*}
因此数列递增且始终小于 $\sqrt2$，极限存在。令极限 $c_\infty$满足同一递推，交叉相乘得 $c_\infty(2c_\infty+3)=3c_\infty+4$，约去共同一次项，得到 $2c_\infty^2=4$。取正根，得到
\[N_\infty=\sqrt2\,mg.\]
总质量趋于无穷却得到有限支持力，是因为上层接近自由落体，传给下层的有效重量递减；它不是静止时无穷总重量的平衡。先用有限层定义极限，也避免直接把无穷质量代进整体 $F=Ma$。
\end{solution}

\begin{question}\begin{example}[线性阻力下返回同一高度的总时间]\label{n4e15}
质量 $m$ 的物体以竖直向上速度 $v_0$ 抛出，运动阻力为 $-kv$，$v$ 为向上正的有符号速度。返回出发高度时向下速率为 $v_f$。求总时间，并解释为何结果可以不用显式求解指数方程。
\fig{n4e15q}{上升和下降的速度符号相反，阻力始终为 $-kv$}{\begin{tikzpicture}
\draw[axis](0,0)--(0,3.2)node[above]{$z$};\fill(0,0)circle(2pt)node[below]{$O$};\draw[flow,blue](-.45,.2)--(-.45,1.6)node[left]{$v_0$};\draw[flow,red](.45,1.4)--(.45,.1)node[right]{$v_f$};\draw[guide](-1,0)--(1,0);\fill(0,2.6)circle(2pt)node[right]{最高点};\node at(0,-.65){运动始终沿同一竖直直线};
\end{tikzpicture}}
\end{example}\end{question}
\begin{solution}\textbf{对有符号速度方程作全程积分}\par\nobreak
图\ref{n4e15q}向上为正，全程方程统一为 $m\dot v=-mg-kv$，不需要在转向点人为改变 $k$ 的符号。由初态到返回时刻 $T$ 积分，速度端点为 $v_0$ 与 $-v_f$：
\[
m\int_{v_0}^{-v_f}\dd w=-mg\int_0^T\dd t-k\int_0^Tv\dd t.
\]
左边为 $m[w]_{v_0}^{-v_f}$，重力项为 $-mg[t]_0^T$，因此
\[m(-v_f-v_0)=-mgT-k\int_0^Tv\dd t.\]


\par\medskip\Needspace{4\baselineskip}\noindent\textbf{用零位移消去阻力总冲量}\par\nobreak
因为返回同一高度，$\int_0^Tv\dd t=\Delta z=0$，所以
\[T=\frac{v_0+v_f}{g}.\]
\fig{n4e15s}{速度图的正、负面积相抵，线性阻力的总冲量因此为零}{\begin{tikzpicture}[x=1.4cm,y=.8cm]
\draw[axis](0,0)--(3.5,0)node[right]{$t$};\draw[axis](0,-1.4)--(0,2.7)node[above]{$v$};\fill[blue!10](0,0)--plot[domain=0:1.098612,samples=40](\x,{3*exp(-\x)-1})--cycle;\fill[red!8](1.098612,0)--plot[domain=1.098612:2.82144,samples=40](\x,{3*exp(-\x)-1})--(2.82144,0)--cycle;\draw[world,domain=0:2.82144,samples=60]plot(\x,{3*exp(-\x)-1});\node[below]at(2.82144,0){$T$};\node[left]at(0,2){$v_0$};\node[right]at(2.82144,-.82){$-v_f$};
\end{tikzpicture}}


\par\medskip\Needspace{4\baselineskip}\noindent\textbf{区分冲量抵消和机械能损失}\par\nobreak
图\ref{n4e15s}说明抵消的是线性阻力的总冲量，不是阻力做功。阻力做功为 $-k\int v^2\dd t<0$，所以 $v_f<v_0$，$T<2v_0/g$。若换成平方阻力，阻力冲量不再正比于位移，上述简化通常失效。给出的 $v_0,v_f,k$ 还须属于同一个可实现过程，它们并非任意独立。
\end{solution}

\begin{question}\begin{example}[竖直弹簧压缩过程的最大速率]\label{n4e16}
质量 $m$ 的物体放在竖直轻弹簧上，弹簧劲度系数 $k$、下端固定，初始原长、物体由静止释放。忽略阻力，求最大速率及其出现位置，检查最大压缩量。
\fig{n4e16q}{向下位移从弹簧原长位置起算}{\begin{tikzpicture}
\draw[world](-1.3,0)--(1.3,0);\node[draw](A)at(0,2.1){$m$};\draw[spring](0,0)--(A.south);\draw[axis](1.5,2.1)--(1.5,-.2)node[below]{$x$};\node[right]at(1.5,2.1){$0$};\draw[guide](-.7,1.2)--(1.5,1.2);\node[below left]at(0,1.2){压缩后位置};
\end{tikzpicture}}
\end{example}\end{question}
\begin{solution}\textbf{用合力变号判断最大速率位置}\par\nobreak
按图\ref{n4e16q}向下取 $x$，合力为 $mg-kx$，加速度在 $x=mg/k$ 变号。物体从零速度向下加速，越过该位置后仍向下但减速，所以最大速率出现在 $x_*=mg/k$。
\fig{n4e16s}{速度平方是位移的开口向下二次函数}{\begin{tikzpicture}
\draw[axis](0,0)--(5,0)node[right]{$x$};\draw[axis](0,0)--(0,2.5)node[above]{$v^2$};\draw[world,domain=0:4,samples=50]plot(\x,{.5*\x*(4-\x)});\draw[guide](2,0)--(2,2);\node[below]at(2,0){$mg/k$};\node[below]at(4,0){$2mg/k$};\node[left]at(0,2){$mg^2/k$};
\end{tikzpicture}}


\par\medskip\Needspace{4\baselineskip}\noindent\textbf{由速度平方曲线求极值和最大压缩量}\par\nobreak
用 $a=(\dd v/\dd x)(\dd x/\dd t)=v\dd v/\dd x$，从初态 $(0,0)$ 积分：
\[
m\int_0^v w\,\dd w=\int_0^x(mg-k\xi)\,\dd\xi,
\qquad
\frac12mv^2=[mg\xi-k\xi^2/2]_0^x=mgx-\frac12kx^2.
\]
所以 $v^2=2gx-kx^2/m$。把 $x_*=mg/k$ 代入，图\ref{n4e16s}顶点处
$v_{\max}^2=2mg^2/k-mg^2/k=mg^2/k$，因此 $v_{\max}=g\sqrt{m/k}$。
令 $v^2=x(2g-kx/m)=0$，除初始根 $x=0$ 外，下一个零速点为 $x_{\max}=2mg/k$。量纲 $g\sqrt{m/k}$ 为速度，且最大速率不在最大压缩处，二者分别由加速度为零和速度为零确定。
\end{solution}

\begin{question}\begin{example}[线性阻力下速度随位移的变化]\label{n4e17}
质量 $m$ 的物体在水平面上以 $v_0>0$ 向右运动，唯一水平作用为线性阻力 $-kv$。从出发位置 $x=0$ 起求 $v(x)$，并说明何时、何处达到零速度。
\fig{n4e17q}{阻力与速度反向，没有恒定库仑摩擦项}{\begin{tikzpicture}
\draw[world](-3,0)--(3,0);\node[draw,anchor=south](A)at(0,0){$m$};\draw[flow](.5,1)--(2,1)node[right]{$v$};\draw[force,red](A.west)--(-2,0 |- A.center)node[left]{$kv$};\draw[axis](-2,-.6)--(2.5,-.6)node[right]{$x$};
\end{tikzpicture}}
\end{example}\end{question}
\begin{solution}\textbf{在正速度区间转换微分变量}\par\nobreak
图\ref{n4e17q}在速度仍正的区间，链式法则给
$\dot v=(\dd v/\dd x)\dot x=v\dd v/\dd x$，故
$mv\dd v/\dd x=-kv$。约去非零的 $v$，并从 $(0,v_0)$ 积分，
\[
\int_{v_0}^{v}\dd w=-\frac km\int_0^x\dd\xi,
\qquad
[w]_{v_0}^{v}=-\frac km[\xi]_0^x,
\]
得到 $v(x)=v_0-kx/m$。
\fig{n4e17s}{有限的极限位移对应无限长的衰减时间}{\begin{tikzpicture}
\draw[axis](0,0)--(4.8,0)node[right]{$x$};\draw[axis](0,0)--(0,2.4)node[above]{$v$};\draw[world](0,2)--(4,0);\draw[guide](4,0)--(4,2);\node[left]at(0,2){$v_0$};\node[below]at(4,0){$mv_0/k$};\node at(2.9,1.8){$t\to\infty$};
\end{tikzpicture}}


\par\medskip\Needspace{4\baselineskip}\noindent\textbf{回到时间方程判断能否有限时刻停车}\par\nobreak
图\ref{n4e17s}的线性关系适用于 $0\le x<mv_0/k$。时间方程为 $\dot v=-(k/m)v$，分离变量并代入初值：
\[
\int_{v_0}^{v}\frac{\dd w}{w}=-\frac km\int_0^t\dd s,
\qquad
\ln v-\ln v_0=-\frac{kt}{m}.
\]
因此 $v(t)=v_0\e^{-kt/m}$。从 $x(0)=0$ 再积分，
\[
x(t)=v_0\int_0^t\e^{-ks/m}\dd s
=v_0\left[-\frac mk\e^{-ks/m}\right]_0^t
=\frac{mv_0}{k}(1-\e^{-kt/m}).
\]
所以任何有限时刻 $v>0$，仅在 $t\to\infty$ 时趋于零，位置趋于 $mv_0/k$。把直线延到负速度，会违反最初约去 $v$ 时的运动区间，也会错误地让纯阻力主动驱动物体反向。
\end{solution}

\begin{question}\begin{example}[从有符号速度测量求合力]\label{n4e18}
质量 $m$ 的质点沿固定直线运动，测得两种过程的有符号速度分别为 $v(t)=v_0+\beta t^3$ 和 $v(t)=v_0\sin(\omega t)$。求各自合力随时间的函数，说明合力图与速度图的关系。
\fig{n4e18q}{速度图提供变化率，不能直接作为力的图像}{\begin{tikzpicture}
\begin{scope}\draw[axis](0,0)--(2.7,0)node[right]{$t$};\draw[axis](0,0)--(0,2.2)node[above]{$v$};\draw[world,domain=0:2,samples=30]plot(\x,{.5+.18*\x*\x*\x});\node[left]at(0,.5){$v_0$};\end{scope}
\begin{scope}[shift={(4,0)}]\draw[axis](0,0)--(3,0)node[right]{$t$};\draw[axis](0,-1.2)--(0,1.5)node[above]{$v$};\draw[world,domain=0:2.7,samples=50]plot(\x,{sin(160*\x)});\end{scope}
\end{tikzpicture}}
\end{example}\end{question}
\begin{solution}\textbf{对多项式速度求导并检查单位}\par\nobreak
图\ref{n4e18q}第一种过程中，$v_0,\beta$均为常量。逐项对时间求导，
\[
 a=\frac{\dd(v_0+\beta t^3)}{\dd t}
 =0+\beta(3t^2)=3\beta t^2,
 \qquad F(t)=ma=3m\beta t^2.
\]
$v_0$只改变初始速度，不改变该过程给定的加速度。若 $\beta$单位为 $\mathrm{m/s^4}$，乘 $mt^2$后力单位正确。
\fig{n4e18s}{正弦速度的合力比速度在相位上提前四分之一周期}{\begin{tikzpicture}
\begin{scope}\draw[axis](0,0)--(2.7,0)node[right]{$t$};\draw[axis](0,0)--(0,2.2)node[above]{$F$};\draw[world,domain=0:2,samples=30]plot(\x,{.45*\x*\x});\node at(1.4,2.4){$3m\beta t^2$};\end{scope}
\begin{scope}[shift={(4,0)}]\draw[axis](0,0)--(3,0)node[right]{$t$};\draw[axis](0,-1.3)--(0,1.7)node[above]{$F$};\draw[world,domain=0:2.7,samples=50]plot(\x,{cos(160*\x)});\node at(1.5,2){$mv_0\omega\cos\omega t$};\end{scope}
\end{tikzpicture}}


\par\medskip\Needspace{4\baselineskip}\noindent\textbf{对正弦速度求导并解释相位关系}\par\nobreak
第二种过程令 $\phi=\omega t$，链式法则给
\[
a=\frac{\dd(v_0\sin\phi)}{\dd\phi}\frac{\dd\phi}{\dd t}
=(v_0\cos\phi)\omega=v_0\omega\cos\omega t.
\]
故图\ref{n4e18s}合力为 $mv_0\omega\cos\omega t$。在速度达到极值时合力为零，在速度过零时合力大小最大。若题目只给没有方向的速率，就不能把正弦负半周直接当作速率；本题明确采用固定轴上的速度分量。所得是合力，不能单凭这些测量进一步唯一分解成驱动力与阻力。
\end{solution}

\section{本讲小结}
多过程动力学的关键是事件、初值和边界。先按真实受力求每段加速度，再连接连续的位置与速度；碰撞、停止和接触变化各用对应条件判断。图像中的斜率和面积应始终同坐标量纲对应，不能用瞬时力或不适当的平均值替代积分。

变力问题可按自变量选择时间、位置或速度积分；重复过程先确定事件顺序，再建立递推。候选极值、无穷极限和停止位置都要配合模型范围解释。结果除了数值，还应交代方向、适用区间、临界之后的分支，以及至少一项独立的物理检查。
\section{练习}
\begin{question}\begin{drill}[平方阻力下的往返时间]\label{n4d1}
物体竖直向上初速 $v_0$，阻力改为 $-cv|v|$，$c>0$。求上升时间、下落到出发点的时间及返回速率，判断总时间是否仍等于 $(v_0+v_f)/g$。可令 $u=\sqrt{mg/c}$。
\fig{n4d1q}{平方阻力在上升与下降时使用不同的速度分支}{\begin{tikzpicture}
\fill(-1.5,0)circle(3pt);\draw[flow,blue](-1.5,.2)--(-1.5,1.5)node[above]{$v>0$};\draw[force,red](-1.5,0)--(-1.5,-1.4)node[below]{$mg+cv^2$};\fill(1.5,0)circle(3pt);\draw[flow,blue](1.5,-.2)--(1.5,-1.5)node[below]{$v<0$};\draw[force,red](1.2,0)--(1.2,-1)node[left]{$mg$};\draw[force](1.8,0)--(1.8,1)node[above]{$cv^2$};
\end{tikzpicture}}\answerlink{n4d1}\end{drill}\end{question}
\begin{question}\begin{drill}[缓慢增加拉力后的滑离时间]\label{n4d2}
上块质量 $m$、长板质量 $2m$，地面光滑，层间最大静摩擦和动摩擦均为 $\mu mg$。两者初静止，拉板的水平力为 $F(t)=\gamma t$，$\gamma>0$。求开始滑动时刻，以及从开始滑动到上块从板左端滑出的时间；初始上块到板左端的相对距离为 $\ell$。
\fig{n4d2q}{板加速更快，上块相对板向左滑动}{\begin{tikzpicture}
\node[draw,minimum width=4cm,minimum height=.6cm]at(0,.3){长板 $2m$};\node[draw,anchor=south]at(.6,.6){$m$};\draw[force](2,.3)--(3.3,.3)node[right]{$\gamma t$};\draw[dim](-2,1.6)--(.6,1.6)node[midway,above]{$\ell$};\draw[world](-2.5,0)--(2.5,0);
\end{tikzpicture}}\answerlink{n4d2}\end{drill}\end{question}
\begin{question}\begin{drill}[不等质量楔块的无穷层极限]\label{n4d3}
沿用例题的光滑 $45^\circ$ 叠层，每层下三角块质量为 $m$，上三角块为 $\eta m$，$\eta>0$。求有限层底部支持系数的递推及无穷层极限，检验 $\eta=1$、$\eta\to0$ 和 $\eta\to\infty$。
\fig{n4d3q}{每层上、下三角块的质量比固定}{\begin{tikzpicture}
\foreach \j in {0,1}{\draw[world](0,{2*\j})rectangle(2,{2*(\j+1)});\draw[world](0,{2*\j})--(2,{2*(\j+1)});\node at(.5,{2*\j+1.5}){$\eta m$};\node at(1.5,{2*\j+.5}){$m$};}\draw[guide](.4,4)--(.4,4.6) (1.6,4)--(1.6,4.6);\draw[force,blue](1,-1)--(1,0)node[right]{$N_n$};
\end{tikzpicture}}\answerlink{n4d3}\end{drill}\end{question}
\begin{question}\begin{drill}[粗糙直槽的最短下滑时间]\label{n4d4}
定点到倾角 $\alpha$ 斜面的法向距离为 $h$，直槽与竖直向下的右偏角限定为 $\beta\ge0$。槽内最大静摩擦与动摩擦均为 $\mu N$，物块从静止向下滑。求最优槽角，说明为何可能出现 $\beta=0$ 的边界最优。
\fig{n4d4q}{重力的法向分量决定槽内摩擦}{\begin{tikzpicture}
\draw[world](-2,0)--(3,2);\draw[world](-1,2.5)--(-.4,.64);\draw[guide](-1,2.5)--(-1,.3);\fill(-.7,1.55)circle(3pt);\draw[force,red](-.7,1.55)--(-.7,.4)node[right]{$mg$};\draw[force,blue](-.7,1.55)--(-1,2.48)node[left]{$\mu N$};\node at(-.8,2){$\beta$};\node at(-1.2,.2){$\alpha$};
\end{tikzpicture}}\answerlink{n4d4}\end{drill}\end{question}

\LectureAnswerSection

\begin{drillsolution}{n4d1}\textbf{分离上升段变量求时间和最高点}\par\nobreak
图\ref{n4d1q}上升段用向上速度 $v>0$，$\dot v=-g(1+v^2/u^2)$。令 $r=v_0/u$，速度由 $v_0$ 降到零，分离变量给
\[
t_\uparrow=\int_0^{v_0}\frac{\dd v}{g(1+v^2/u^2)}
=\frac ug\int_0^r\frac{\dd y}{1+y^2}
=\frac ug[\arctan y]_0^r=\frac ug\arctan r.
\]
这里换元 $y=v/u$，$\dd v=u\dd y$，速度端点 $0,v_0$ 对应 $y=0,r$；原函数来自 $(\arctan y)'=1/(1+y^2)$。
同一上升段用 $v\dd v/\dd z=\dot v$，得
\[
H=\int_0^{v_0}\frac{v\dd v}{g(1+v^2/u^2)}
=\frac{u^2}{2g}\int_1^{1+r^2}\frac{\dd q}{q}
=\frac{u^2}{2g}[\ln q]_1^{1+r^2}
=\frac{u^2}{2g}\ln(1+r^2).
\]
第二个积分换元 $q=1+v^2/u^2$，所以 $\dd q=2v\dd v/u^2$，上下限变为 $1$ 与 $1+r^2$。
两个积分的分母相同，分子多出的 $v$ 来自位移微元 $\dd z=v\dd t$。
\fig{n4d1s}{下落速率渐近终端值，返回速率由有限高度确定}{\begin{tikzpicture}
\draw[axis](0,0)--(4.8,0)node[right]{$t_\downarrow$};\draw[axis](0,0)--(0,2.5)node[above]{$w$};\draw[guide](0,2)--(4.4,2)node[right]{$u$};\draw[world,domain=0:4.2,samples=50]plot(\x,{2*(exp(2*\x)-1)/(exp(2*\x)+1)});\draw[guide](2.2,0)--(2.2,1.95);\node[below]at(2.2,0){返回时刻};
\end{tikzpicture}}



\par\medskip\Needspace{4\baselineskip}\noindent\textbf{从最高点重新积分下降运动}\par\nobreak
从最高点另起计时，记下降用时为 $\tau$，向下速率为 $w$、下降距离为 $s$。
阻力向上，故 $\dot w=g(1-w^2/u^2)$，初始条件为 $w(0)=s(0)=0$。
在 $0\le w<u$ 的实际分支上分离变量，令 $z=w/u$，得
\[
 \frac{g\tau}{u}=\int_0^{w/u}\frac{\dd z}{1-z^2}
 =\frac12[-\ln(1-z)+\ln(1+z)]_0^{w/u}
 =\frac12\ln\frac{1+w/u}{1-w/u}.
\]
这里用了 $1/(1-z^2)=[1/(1-z)+1/(1+z)]/2$。
将对数式指数化并解出 $w$，得到
\[
 w=u\frac{\e^{2g\tau/u}-1}{\e^{2g\tau/u}+1}
   =u\tanh\frac{g\tau}{u}.
\]
定义 $\cosh x=(\e^x+\e^{-x})/2$、
$\tanh x=(\e^x-\e^{-x})/(\e^x+\e^{-x})$。链式求导给
$\dd\ln\cosh x/\dd x=(\e^x-\e^{-x})/(2\cosh x)=\tanh x$。
位移积分中换元 $q=g\tau'/u$，有 $\dd\tau'=(u/g)\dd q$，
原上下限 $0,\tau$ 对应 $0,g\tau/u$，所以
\[
 s=\int_0^\tau w(\tau')\dd\tau'
   =\frac{u^2}{g}\int_0^{g\tau/u}\tanh q\,\dd q
   =\frac{u^2}{g}[\ln\cosh q]_0^{g\tau/u}
   =\frac{u^2}{g}\ln\cosh\frac{g\tau}{u}.
\]
图\ref{n4d1s}中的 $w$ 从零单调增大，却始终低于终端速率 $u$；
因此不能把下降段也写成上升时含 $1+w^2/u^2$ 的分母。

\par\medskip\Needspace{4\baselineskip}\noindent\textbf{由返回高度确定下降时间与末速}\par\nobreak
令 $s=H$，并记 $x=gt_\downarrow/u\ge0$，则
\[
 \ln\cosh x=\frac12\ln(1+r^2),\qquad
 \cosh x=\sqrt{1+r^2}.
\]
把 $\cosh x=(\e^x+\e^{-x})/2$ 代入，令 $Y=\e^x\ge1$，可得
$Y^2-2\sqrt{1+r^2}Y+1=0$。
两个根互为倒数；下降时间非负，故应取
$Y=\sqrt{1+r^2}+r$。再由 $\tanh x=\sqrt{\cosh^2x-1}/\cosh x$
求返回速率，图\ref{n4d1s}的返回状态为
\[
 t_\downarrow=\frac ug\ln\left(r+\sqrt{1+r^2}\right),\qquad
 v_f=\frac{v_0}{\sqrt{1+r^2}}.
\]


\par\medskip\Needspace{4\baselineskip}\noindent\textbf{整理总时间并检查弱阻力极限}\par\nobreak
下降时间只从最高点算起；从发射到返回的总时间是
\[
 T=t_\uparrow+t_\downarrow
 =\frac ug\left[\arctan r+\ln\left(r+\sqrt{1+r^2}\right)\right].
\]
它一般不等于 $(v_0+v_f)/g$，因为平方阻力的冲量不能用线性阻力中的 $-k\Delta z$ 表示。
$c\to0$、$r\to0$ 时，$\arctan r\simeq r$、
$\ln(r+\sqrt{1+r^2})\simeq r$，两段时间都趋于 $v_0/g$，且 $v_f\to v_0$，回到真空抛体。
\end{drillsolution}

\begin{drillsolution}{n4d2}\textbf{以静摩擦上限确定开始滑动的时刻}\par\nobreak
图\ref{n4d2q}共同运动时 $a=F/(3m)$，上块所需摩擦 $f=ma=F/3$。达到 $\mu mg$ 时 $F_c=3\mu mg$，所以 $t_c=3\mu mg/\gamma$。
\fig{n4d2s}{滑动后的相对位移从零以三次方增长}{\begin{tikzpicture}
\draw[axis](0,0)--(4.7,0)node[right]{$\tau=t-t_c$};\draw[axis](0,0)--(0,2.7)node[above]{$s_{\rm rel}$};\draw[world,domain=0:3.5,samples=40]plot(\x,{.05*\x*\x*\x});\draw[guide](0,2.14)--(3.5,2.14)--(3.5,0);\node[left]at(0,2.14){$\ell$};\node[below]at(3.5,0){$\tau_e$};
\end{tikzpicture}}


\par\medskip\Needspace{4\baselineskip}\noindent\textbf{在局部时间内积分相对加速度}\par\nobreak
滑动后上块加速度 $a_m=\mu g$，板加速度 $a_P=[\gamma t-\mu mg]/(2m)$。令 $\tau=t-t_c$，两者之差 $a_P-a_m=\gamma\tau/(2m)$。临界时相对速度为零，以从临界位置新增的滑动量定义 $s_{\rm rel}(0)=0$，逐次积分：
\begin{align*}
v_P-v_m&=\int_0^\tau\frac{\gamma s}{2m}\dd s
=\left[\frac{\gamma s^2}{4m}\right]_0^\tau
=\frac{\gamma\tau^2}{4m},\\
s_{\rm rel}&=\int_0^\tau\frac{\gamma s^2}{4m}\dd s
=\left[\frac{\gamma s^3}{12m}\right]_0^\tau
=\frac{\gamma\tau^3}{12m}.
\end{align*}



\par\medskip\Needspace{4\baselineskip}\noindent\textbf{用板端边界求总离板时刻}\par\nobreak
图\ref{n4d2s}取 $s_{\rm rel}=\ell$，得滑动后到左端时间 $\tau_e=(12m\ell/\gamma)^{1/3}$，总时刻 $t_e=t_c+\tau_e$。整个滑动阶段相对速度为正且递增，因此摩擦方向一致；如果只在临界时把恒定加速度带入位移公式，会遗漏持续增加的驱动力。
\end{drillsolution}

\begin{drillsolution}{n4d3}\textbf{由质量比与斜面约束求接触力分量}\par\nobreak
定义图\ref{n4d3q}的 $N_n=c_nmg$，$c_0=0$。设上块向下加速度 $a$，接触斜面力水平与竖直分量大小相同，记为 $P$。两块水平加速度大小分别为 $P/m$、$P/(\eta m)$，$45^\circ$ 约束给两者之和等于 $a$，所以 $P=bma$，$b=\eta/(1+\eta)$。
\fig{n4d3s}{质量比改变水平分配，但上层仍由有效重力缩放}{\begin{tikzpicture}
\draw[world](0,0)rectangle(2,2);\draw[world](0,0)--(2,2);\draw[flow](.5,1.5)--(-.8,1.5)node[left]{$P/(\eta m)$};\draw[flow](1.5,.5)--(2.8,.5)node[right]{$P/m$};\draw[force,red](1,2.9)--(1,2)node[midway,right]{$c_{n-1}m(g-a)$};\node at(1,-.7){$N_n=mg+P$};
\end{tikzpicture}}



\par\medskip\Needspace{4\baselineskip}\noindent\textbf{使用上层有效重力建立递推}\par\nobreak
图\ref{n4d3s}上块竖直式为 $\eta ma=\eta mg+c_{n-1}m(g-a)-P$，得到 $a/g=(\eta+c_{n-1})/(\eta+c_{n-1}+b)$。因此
\[c_n=1+b\frac{\eta+c_{n-1}}{\eta+c_{n-1}+b}.\]


\par\medskip\Needspace{4\baselineskip}\noindent\textbf{证明极限存在并核查质量比极限}\par\nobreak
记递推映射为 $f(c)=1+b(\eta+c)/(\eta+c+b)$。商法则给
\[
f'(c)=b\frac{(\eta+c+b)-(\eta+c)}{(\eta+c+b)^2}
=\frac{b^2}{(\eta+c+b)^2}>0.
\]
因 $\eta>0,c\ge0$，此导数小于 $1$。而 $c_1>c_0=0$，递增映射使 $c_{n+1}>c_n$；又有 $f(c)<1+b$，故有限层序列递增有界，极限存在。
令 $c$ 为极限，固定点式给
$(c-1)(\eta+c+b)=b(\eta+c)$。整理成
$(c-1)(c+\eta)=b(\eta+1)=\eta$，
即 $c^2+(\eta-1)c-2\eta=0$。舍去负根，所以
\[N_\infty=\frac{1-\eta+\sqrt{\eta^2+6\eta+1}}2\,mg.\]
$\eta=1$ 给 $\sqrt2mg$；$\eta\to0$ 直接代入给 $mg$。为计算 $\eta\to\infty$，将正根有理化，
\[
c=\frac{4\eta}{\sqrt{\eta^2+6\eta+1}+\eta-1}
=\frac4{\sqrt{1+6/\eta+1/\eta^2}+1-1/\eta}\longrightarrow2,
\]
所以支持力趋于 $2mg$。后两项是先取无限层极限再取质量比极限，说明上块很重时也不能把底部支持直接按无限总重力计算。
\end{drillsolution}

\begin{drillsolution}{n4d4}\textbf{隔离槽内物块并结合路径几何}\par\nobreak
图\ref{n4d4q}直槽不弯曲，物块法向加速度为零，故 $N=mg\sin\beta$，摩擦沿槽向上。沿槽加速度 $a_s=g(\cos\beta-\mu\sin\beta)$，要求括号为正。槽长仍为 $s=h/\cos(\alpha-\beta)$，所以
\[t^2=\frac{2h}{g(\cos\beta-\mu\sin\beta)\cos(\alpha-\beta)}.\]
\fig{n4deep:roughslot}{沿槽重力分量驱动下滑，法向分量决定摩擦上限}{\begin{tikzpicture}
\coordinate(P)at(0,0);\draw[world](-1,2)--(1,-2);\fill(P)circle(2pt);\draw[guide](0,0)--(0,-2.3);\draw[force,red](P)--(0,-1.8)node[below]{$mg$};\draw[force,blue](P)--(1.2,.6)node[right]{$N$};\draw[force](P)--(-.6,1.2)node[left]{$\mu N$};\draw[flow](.8,-.6)--(1.4,-1.8)node[right]{$a_s$};\node at(.23,-.9){$\beta$};\node at(2.1,1.5){$N=mg\sin\beta$};
\end{tikzpicture}}
图\ref{n4deep:roughslot}中槽与竖直夹角为 $\beta$，支持力与槽垂直；于是重力沿槽分量为 $mg\cos\beta$，法向分量为 $mg\sin\beta$。摩擦随 $\beta$ 增大而增大，这解释了下面最优方向为何向竖直靠拢。

\fig{n4d4s}{摩擦把内部最优方向推向竖直，超过后由边界接替}{\begin{tikzpicture}
\draw[axis](0,0)--(4.8,0)node[right]{$\delta=\arctan\mu$};\draw[axis](0,0)--(0,2.2)node[above]{$\beta_*$};\draw[world](0,1.6)--(3,0)--(4.3,0);\node[left]at(0,1.6){$\alpha/2$};\node[below]at(3,0){$\alpha$};
\end{tikzpicture}}



\par\medskip\Needspace{4\baselineskip}\noindent\textbf{求内部极值并检查角度边界}\par\nobreak
令 $\delta=\arctan\mu$，$\cos\beta-\mu\sin\beta=\sqrt{1+\mu^2}\cos(\beta+\delta)$。分母除去正常数后为 $[\cos(\alpha+\delta)+\cos(2\beta+\delta-\alpha)]/2$。当 $\alpha\ge\delta$，内部最大值在 $\beta_*=(\alpha-\delta)/2$；当 $\alpha<\delta$，这个角为负，超出题设范围，图\ref{n4d4s}最优改为边界 $\beta_*=0$。



\par\medskip\Needspace{4\baselineskip}\noindent\textbf{验证运动条件及光滑极限}\par\nobreak
内部解满足 $\cos(\beta_*+\delta)=\cos[(\alpha+\delta)/2]>0$，故确能下滑并到达斜面。边界解竖直槽的法向力为零，也就没有该模型中的摩擦，时间为 $\sqrt{2h/(g\cos\alpha)}$。$\mu\to0$ 时回到光滑直槽的 $\alpha/2$，验证参数极限。
\end{drillsolution}

\end{document}

